%% ============================================================
%%  SUPPLEMENTARY MATERIAL B
%%  HypatiaX: Noise Robustness, Sample Complexity, and Convergence
%%  of Hybrid LLM–Neural Symbolic Regression on a Custom 30-Equation Benchmark
%%
%%  Companion to: jmlr_paper_main.tex
%%  Sources: six-method noiseless comparison (clean report)  +
%%           noise/sample-complexity sweep (v2, 16 March 2026)
%%  Rewritten & restructured: March 2026
%% ============================================================

\documentclass[twoside,11pt]{article}

%% ── JMLR house style ─────────────────────────────────────────────────────────
\usepackage{jmlr2e}

%% ── Core packages ────────────────────────────────────────────────────────────
\usepackage{amsmath,amssymb}
\usepackage{booktabs}
\usepackage{multirow}
\usepackage{makecell}
\usepackage{xcolor}
\usepackage{colortbl}
\usepackage{array}
\usepackage{longtable}% generated per-case tables (llm_detailed, defi_detailed)
\usepackage{caption}
\usepackage{subcaption}
\usepackage[expansion=false]{microtype}
\usepackage{rotating}
\usepackage{algorithm}
\usepackage{algorithmic}
\usepackage{graphicx}
\usepackage{float}
\usepackage{pgfplots}
\pgfplotsset{compat=1.18}
\usepackage{cleveref}
% Cross-document references to the main paper are hard-coded (Table 17,
% Sections 10.1 and 10.9; not to be confused with this document's own Table 17) because this file does not load xr. Re-check them if
% the main paper's numbering changes.
\usepackage{verbatim}
\usepackage{listings}

%% ── Code listings ────────────────────────────────────────────────────────────
\definecolor{codeblue}{RGB}{0,90,170}
\definecolor{codegray}{RGB}{100,100,100}
\definecolor{codegreen}{RGB}{0,128,64}
\definecolor{codered}{RGB}{180,30,30}
\definecolor{backgray}{RGB}{248,248,248}

\lstset{
  language=Python,
  basicstyle=\ttfamily\small,
  keywordstyle=\color{codeblue}\bfseries,
  commentstyle=\color{codegray}\itshape,
  stringstyle=\color{codegreen},
  backgroundcolor=\color{backgray},
  frame=single,
  framerule=0.4pt,
  rulecolor=\color{black!20},
  breaklines=true,
  breakatwhitespace=true,
  numbers=left,
  numberstyle=\tiny\color{codegray},
  numbersep=8pt,
  xleftmargin=12pt,
  tabsize=4,
  showstringspaces=false,
  captionpos=b
}

%% ── Colours ──────────────────────────────────────────────────────────────────
\definecolor{passgreen}{RGB}{22,163,74}
\definecolor{passgreenBg}{RGB}{220,252,231}
\definecolor{failred}{RGB}{220,38,38}
\definecolor{failredBg}{RGB}{254,226,226}
\definecolor{warnAmber}{RGB}{217,119,6}
\definecolor{warnAmberBg}{RGB}{254,243,199}
\definecolor{headnavyBg}{RGB}{219,234,254}
\definecolor{m3blue}{RGB}{31,119,180}
\definecolor{m4red}{RGB}{214,39,40}
\definecolor{tabhead}{RGB}{220,230,245}

%% ── Column types ─────────────────────────────────────────────────────────────
\newcolumntype{C}[1]{>{\centering\arraybackslash}p{#1}}
\newcolumntype{L}[1]{>{\raggedright\arraybackslash}p{#1}}
\newcolumntype{R}[1]{>{\raggedleft\arraybackslash}p{#1}}

%% ── Method macros ────────────────────────────────────────────────────────────
\newcommand{\PureLLM}{\textsc{PureLLM}}
\newcommand{\INN}{\textsc{ImpNN}}
\newcommand{\EHD}{\textsc{EHSDeFi}}
\newcommand{\HSL}{\textsc{HLLMNN}}
\newcommand{\SEL}{\textsc{SymLLM}}
\newcommand{\HDS}{\textsc{HDSv50\_2}}
\newcommand{\Mthree}{{\color{m3blue}\textbf{\EHD}}}
\newcommand{\Mfour}{{\color{m4red}\textbf{\HSL}}}

%% ── Metric shorthand ─────────────────────────────────────────────────────────
\newcommand{\Rsq}{R^{2}}
\newcommand{\rmse}{\ensuremath{\mathrm{RMSE}}}
\newcommand{\nrmse}{\ensuremath{\widetilde{\mathrm{RMSE}}}}
\newcommand{\warn}{{\color{warnAmber}\textbf{!}}}

%% ── Theorem-like environments ────────────────────────────────────────────────
%% NOTE: amsthm is deliberately NOT loaded and `example` is NOT
%% redeclared here. jmlr2e.sty already defines proof/theorem/lemma/
%% proposition/remark/corollary/definition/example unconditionally
%% via plain-LaTeX \newtheorem/\newenvironment. Loading amsthm here
%% interferes with jmlr2e's optional-argument counter machinery
%% (symptom: "No counter 'Definition' defined" elsewhere in the
%% submission), and redeclaring `example` collides outright
%% ("Command \theorem already defined"-class error). Neither
%% needed — jmlr2e's `example` environment (shared, continuous
%% numbering, not per-section) already supports the
%% \begin{example}[Title] usage below.

%% ============================================================
\usepackage{xr-hyper}
\externaldocument{jmlr_paper_main_submission}
\begin{document}

\title{%
  Supplementary Material B\\[6pt]
  HypatiaX: Noise Robustness, Sample Complexity, and Convergence\\
  of Hybrid LLM--Neural Symbolic Regression on a Custom 30-Equation Benchmark
}
\author{Ruperto Pedro Bonet Chaple\\
  \texttt{ruperto.bonet@modelphysmat.com}\\
  Independent Researcher, London, United Kingdom
}
\editor{}
\maketitle

\begin{abstract}
We present an extended evaluation of the HypatiaX symbolic regression
framework on a custom 30-equation benchmark, integrating a six-method
noiseless comparison with a comprehensive noise robustness sweep
($\sigma \in \{0, 0.05, 0.1, 0.5, 1\}\%$) and a
sample complexity sweep
($n \in \{50, 100, 200, 500, 750, 1000\}$), focusing on the two
top-performing core methods: \EHD{} (M3) and \HSL{} (M4).

Following correction of a measurement harness bug in \HSL{} (see
Section~\ref{sec:bugfix}), both methods achieve median $\Rsq = 1.000000$
across all tested conditions.
\EHD{} achieves 100\% recovery at all \emph{noisy} conditions ($\sigma > 0$)
but only 76.7\% (23/30) at $\sigma=0\%$ under the strict noiseless
threshold ($\Rsq \geq 0.999999$; \Cref{tab:rr_noise}).

\HSL{} achieves 100\% recovery at all \emph{noisy} conditions ($\sigma > 0$)
and 100\% (30/30 equations, v2-corrected) under the strict noiseless protocol
($\Rsq > 0.9999$). \HDS{} achieves 70.0\% (21/30 equations) under the same
noiseless protocol.

Both reach median $\Rsq=1.000000$ at $n = 50$ samples (noiseless sample-complexity
sweep; noisy $n>50$ runs are not available), though \EHD's
recovery rate dips to 96.7\% at larger $n$ (\Cref{tab:sc_metrics}).
\HSL{} is $\sim$1.1--1.2$\times$
faster at every noise level.

On this 30-equation custom benchmark (\textbf{not} a Feynman subset ---
see \Cref{tab:overall}), \EHD{} achieves 76.7\% recovery and \HDS{} achieves
70.0\%.
We additionally correct three methodological issues:
undisclosed hardcoded \PureLLM{} results, scale-biased raw \rmse,
and mismatched pass-rate thresholds.
\end{abstract}


\begin{keywords}
symbolic regression, scientific discovery, large language models,
neural networks, equation-recovery benchmark, noise robustness, sample complexity,
hybrid systems, reproducibility
\end{keywords}

%% ============================================================
\section{Introduction}
\label{suppb:intro}

This document is Supplementary Material~B to the main JMLR submission
\emph{HypatiaX: A Hybrid Symbolic-Neural Framework for Extrapolation-Reliable
Analytical Discovery} (\texttt{jmlr\_paper\_main.tex}).

Symbolic regression (SR) aims to discover closed-form analytical expressions
from data without prior knowledge of the functional form.  The Feynman SR
Benchmark~\citep{udrescu2020ai} provides a canonical evaluation suite of
100 physics equations (the standard set used under SRBench).

This supplementary integrates two complementary evaluations:
\begin{enumerate}
  \item A \textbf{six-method noiseless comparison} (corrected from the
    companion report). \textbf{[Note:
    this evaluation uses a custom 30-equation benchmark, not a subset of the
    Feynman suite described above. Its 11 domains are biology, chemistry,
    electrochemistry, electromagnetism, electrostatics, magnetism, mechanics,
    optics, probability, quantum, and thermodynamics; several of these
    (biology, chemistry, electrochemistry, probability) are outside the
    Feynman benchmark's scope.  The domain list above is taken from the source files. See
    \Cref{sec:noiseless} and the source-file note on \Cref{tab:overall}.]}
  \item A \textbf{sweep evaluation} 
    covering \EHD{} (M3) and \HSL{} (M4) across noise levels and sample
    sizes, with convergence, win-rate, and architecture analysis.
\end{enumerate}

%% ============================================================
\section{Related Work}
\label{suppb:related}

SR has evolved from genetic programming~\citep{schmidt2009distilling} through
physics-inspired decomposition~\citep{udrescu2020ai} to transformer-based
generation~\citep{biggio2021neural,shojaee2023transformer} and deep RL over
expression trees~\citep{petersen2021deep}.

Benchmark integrity has received increasing attention.
\citet{memorization2024audit} show that LLMs can reproduce iconic equations
from training data rather than inferring them from observations, inflating
reported pass rates.  The hardcoding audit applied to \PureLLM{} in
Section~\ref{sec:integrity} follows this line of work.

A methodological concern specific to tiny-scale physics equations is that
benchmark harnesses using absolute sum-of-squares thresholds silently
misclassify correct formulas.  The fix described in Section~\ref{sec:bugfix}
is a general contribution to SR benchmarking practice.

%% ============================================================
\section{Methods}
\label{sec:methods}

\subsection{Six-Method Benchmark Suite}
\label{sec:methods:sixmethod}

\begin{table}[ht]
\caption{Methods under evaluation: six-method noiseless benchmark.}
\label{tab:methods}
\centering
\small
\renewcommand{\arraystretch}{1.3}
\begin{tabular}{L{2.1cm} L{1.3cm} L{8.0cm}}
\toprule
\textbf{Method} & \textbf{Type} & \textbf{Description} \\
\midrule
\rowcolor{warnAmberBg}
\PureLLM & Core &
  Zero-shot LLM formula inference.  No numerical fitting.
  \textbf{\warn~11/30 results flagged hardcoded (Section~\ref{sec:integrity}).} \\
\EHD\,(M3) & Core &
  Ensemble-first system (140/150 = 93.3\% of routing decisions; \Cref{tab:arch}) combining LLM-guided generation,
  NN ensemble scoring, and DeepFit optimisation. \\
\HSL\,(M4) & Core &
  LLM-first + NN residual correction (123/150 = 82.0\% LLM-first routing decisions; \Cref{tab:arch});
  $\sim$1.1--1.2$\times$ faster than M3 (\Cref{tab:time_noise}). \\
\INN & Core &
  Pure neural baseline ($1\!\to\!128\!\to\!64\!\to\!32\!\to\!1$, 3 seeds). \\
\SEL & Tools &
  LLM-guided symbolic search with CAS verification. \\
\HDS & Tools &
  HybridDiscoverySystem~v50\_2; 21/30 pass at 75.3\,s mean runtime in the noiseless protocol. \\
\bottomrule
\end{tabular}
\end{table}

\subsection{Sweep Evaluation Design (M3 vs.\ M4)}
\label{sec:methods:sweep}

\begin{table}[ht]
\centering
\caption{Sweep design.}
\label{tab:sweeps}
\begin{tabular}{L{2.5cm} L{4.3cm} L{2.5cm} C{2.0cm}}
\toprule
\textbf{Sweep} & \textbf{Axis varied} & \textbf{Fixed} & \textbf{Conditions}\\
\midrule
Noise       & $\sigma \in \{0, 0.05, 0.1, 0.5, 1\}\%$\textsuperscript{*} & $n=200$\textsuperscript{$\dagger$} & 5 levels \\
Sample size & $n \in \{50,100,200,500,750,1000\}$  & $\sigma=5\%$ & 6 sizes  \\
\bottomrule
\multicolumn{4}{l}{\footnotesize\textsuperscript{*}Noise levels as recorded in the sweep source files; see \Cref{tab:r2_noise}.}\\
\end{tabular}
\end{table}

\noindent\textbf{Hyperparameters:} NN seeds = 5, random seed 42,
method timeout = 900\,s, PySR timeout = 1100\,s,
threshold (noise sweep) = 0.999999 at $\sigma=0$, $0.995$ at
$\sigma\in\{0.05,0.1,0.5\}\%$ and $0.99$ at $\sigma=1\%$
(sample-complexity sweep, noiseless: $0.999$ at $n=50$, $0.9995$ at
$n=100$, $0.9999$ at $n\geq200$).
All 30 equations completed with zero timeouts across all 11 conditions.

\subsection{Integrity Annotation: Hardcoded Results}
\label{sec:integrity}

Following~\citet{memorization2024audit}, each \PureLLM{} noiseless row is audited for whether the language model
was called. The \texttt{is\_hardcoded} flag is set when the baseline returns
\texttt{method="pure\_llm\_hardcoded"}, i.e.\ when \texttt{\_try\_hardcoded\_formula} in
\texttt{baseline\_pure\_llm\_defi\_discovery.py} returns a formula without any LLM call. It is a record of
that code path, not a statistical memorisation test. Two kinds of bypass exist: eight equations return the
benchmark's own ground-truth expression (GT in Table~\ref{tab:hardcoded}), and three use a fixed functional form whose parameters are fitted to the
data by least squares (Fit). A twelfth branch (Arrhenius) did not fire. We call the remaining 19 rows \emph{live}.
The same 11 equations are identified in the eleven per-domain \texttt{protocol\_core\_noiseless\_20260812\_*.json} files (by flag), in the consolidated \texttt{benchmark\_results.json} and in \texttt{benchmark\_results\_pca\_4060.json} (by timing, 0.5--1.6\,ms, since those two flat files carry no flag), using \texttt{recount\_hardcoded.py}. The \texttt{170344} file was not part of the recount.
11 of 30 \PureLLM{} noiseless rows bypassed the model (Table~\ref{tab:hardcoded}).

\begin{table}[ht]
\centering
\caption{Equations for which \PureLLM{} bypassed the language model (flagged \texttt{is\_hardcoded=True}).
  GT: the benchmark's ground-truth expression is returned; Fit: fixed form with least-squares parameters fitted to the data.
  $\Rsq$ as recorded in the recount; every bypassed row passes.}
\label{tab:hardcoded}
\small
\begin{tabular}{L{5.4cm} L{3.0cm} C{1.0cm} C{1.7cm}}
\toprule
\textbf{Equation} & \textbf{Domain} & \textbf{Type} & \textbf{$\Rsq$ (NL)} \\
\midrule
Michaelis--Menten enzyme kinetics       & Biology          & Fit & 1.0000000 \\
Logistic growth rate                    & Biology          & Fit & 1.0000000 \\
Allometric scaling law                  & Biology          & Fit & 1.0000000 \\
Nernst equation (electrode potential)   & Electrochemistry & GT  & 1.0000000 \\
Clausius--Mossotti effective field      & Electrostatics   & GT  & 1.0000000 \\
Lorentz force: $F = qvB$                & Magnetism        & GT  & 1.0000000 \\
Gaussian/normal distribution PDF        & Probability      & GT  & 1.0000000 \\
Zeeman energy (electron spin)           & Quantum          & GT  & 1.0000000 \\
Bose--Einstein occupation number        & Quantum          & GT  & 1.0000000 \\
Planck blackbody spectral radiance      & Thermodynamics   & GT  & 1.0000000 \\
Fourier's law of heat conduction        & Thermodynamics   & GT  & 1.0000000 \\
\bottomrule
\end{tabular}
\end{table}

\subsection{Evaluation Metrics}
\label{sec:methods:metrics}

\paragraph{Coefficient of determination ($\Rsq$).}
$\Rsq = 1 - \mathrm{SS}_{\mathrm{res}} / \mathrm{SS}_{\mathrm{tot}}$.

\paragraph{Normalised RMSE ($\nrmse$).}
$\nrmse = \rmse/\sigma_y = \sqrt{1 - \Rsq}$.
Raw \rmse{} is dominated by large-scale equations; median \nrmse{} or $\Rsq$
should be used for cross-equation comparison.

\paragraph{Recovery rate.}
Fraction of equations achieving $\Rsq \geq $ threshold ($\geq 0.9999$
noiseless; $\geq 0.995$ noisy).  Catastrophic = $\Rsq < 0.90$.

%% ============================================================
\section{Algorithm}
\label{sec:algorithm}

\begin{algorithm}[ht]
\caption{HypatiaX hybrid symbolic discovery (M3/M4 core loop)}
\label{alg:hds}
\begin{algorithmic}[1]
\REQUIRE Dataset $\mathcal{D}$, LLM proposer $\mathcal{M}$,
         NN ensemble $f_\theta$, threshold $\tau$
\ENSURE  Discovered expression $\hat{e}$
\STATE $d \leftarrow \textsc{Describe}(\mathcal{D})$
\STATE $\{e_1,\dots,e_K\} \leftarrow \mathcal{M}(d, \mathcal{D})$
\FOR{each candidate $e_k$}
  \STATE Fit $\theta^* \leftarrow \arg\min_\theta \sum_i (f_\theta(e_k,x_i)-y_i)^2$
  \STATE Compute $\Rsq_k$ with relative threshold (Section~\ref{sec:bugfix})
  \IF{$e_k$ invalid or numerically unstable}
    \STATE \textbf{discard}
  \ENDIF
\ENDFOR
\STATE \textbf{M3 path:} $\hat{e} \leftarrow \arg\max_{e_k} \Rsq_k$
\STATE \textbf{M4 path:} if $\Rsq_{\text{LLM}} \geq \tau$, return LLM formula;
       else apply NN residual correction
\RETURN $\hat{e}$
\end{algorithmic}
\end{algorithm}

%% ============================================================
\section{Six-Method Noiseless Results}
\label{sec:noiseless}

\begin{table}[ht]
\caption{Six-method aggregate performance, noiseless protocol
  ($\Rsq \geq 0.999999$, 30 equations).
  \warn~\PureLLM{} pass rate includes 11/30 hardcoded results.
  Source: \texttt{protocol\_core\_noiseless\_20260812\_170344.json} for all rows except the \PureLLM{} pass count. The \PureLLM{} row is a recount over the eleven per-domain \texttt{protocol\_core\_noiseless\_20260812\_*.json} files: 29/30 over all rows, of which 11 are bypassed rows (\Cref{tab:hardcoded}) that all pass, and 18/19 on the \emph{live} rows where the model was called; the 170344 file reports 30/30 (the one-test difference is Curie's law). This is a custom benchmark, not the Feynman-30 suite (\Cref{suppb:intro}); its 11 domains are biology, chemistry, electrochemistry, electromagnetism, electrostatics, magnetism, mechanics, optics, probability, quantum, and thermodynamics.}
\label{tab:overall}
\centering
\small
\renewcommand{\arraystretch}{1.3}
\begin{tabular}{L{2.4cm} C{1.0cm}
                C{1.8cm} C{1.3cm} C{1.4cm}}
\toprule
\textbf{Method} & \textbf{Type} &
  \makecell{\textbf{Pass Rate}\\$(\Rsq \geq 0.999999)$} &
  \makecell{\textbf{Mean} $\Rsq$\\(valid tests)} &
  \makecell{\textbf{Avg}\\Runtime} \\
\midrule
\rowcolor{warnAmberBg}
\PureLLM\,\warn$^\dagger$ & Core  & 29/30 (96.7\%)\warn & --- &   4.7\,s \\
\rowcolor{passgreenBg}
\EHD\,(M3)      & Core  & 23/30 (76.7\%)      & 0.999998 &   7.6\,s \\
\HDS            & Tools & 21/30 (70.0\%)      & 0.999473 &  75.3\,s \\
\HSL\,(M4)      & Core  & 30/30 (100.0\%)     & 1.000000  &   6.9\,s \\
\SEL            & Tools & 28/30 (93.3\%)      & 0.999711 &  39.1\,s \\
\rowcolor{failredBg}
\INN            & Core  & 0/30 (0.0\%)        & 0.999222 &   2.5\,s \\
\bottomrule
\multicolumn{5}{l}{\footnotesize $^\dagger$\,Recount: 29/30 over all rows, 18/19 live rows, 11/11 bypassed rows. Mean $\Rsq$
  and runtime were not recomputed; the 4.7\,s average includes bypassed rows, which take under 2\,ms.}\\
\end{tabular}
\end{table}

\subsection{Scale-Invariant RMSE}
\label{sec:noiseless:nrmse}

After replacing raw \rmse{} with $\nrmse = \sqrt{1-\Rsq}$, four of six methods
cluster in mean $\nrmse \in [0.000, 0.021]$; \PureLLM{} is the outlier on mean
($\nrmse_{\text{mean}}=0.1042$)   \HDS{} has raw \rmse{} values
of \texttt{nan}/sentinel-overflow in the source data for
specific equations (see table note), so its mean raw \rmse{} over all
30 equations is not meaningful; \nrmse{} remains well-defined for every
equation with a finite $\Rsq$ since it is derived directly from $\Rsq$.

\begin{table}[ht]
\centering
\caption{Raw RMSE vs.\ normalised RMSE ($\nrmse=\sqrt{1-\Rsq}$), 30 equations.
Regenerated from the eleven per-domain \texttt{protocol\_core\_noiseless\_20260812\_*.json} files (not the 170344 file that supplies five rows of \Cref{tab:overall}; see its tablenote).}
\label{tab:nrmse}
\small
\begin{tabular}{L{3.0cm} C{2.6cm} C{2.0cm} C{2.0cm}}
\toprule
\textbf{Method} & \textbf{Mean RMSE} & \textbf{Mean NRMSE} & \textbf{Max NRMSE} \\
\midrule
\PureLLM (core) & --- & 0.1042 & --- \\
\EHD     (core) & 0.5672        & 0.0037 & 0.0735 \\
\HDS     (tools)& 2{,}740.0$^\ddagger$ & 0.0199 & 0.1675 \\
\SEL     (tools)& 106.7          & 0.0099 & 0.2091 \\
\INN     (core) & 361.4          & 0.0209 & 0.2419 \\
\HSL     (core) & 0.0002         & 0.0002 & 0.0038 \\
\bottomrule
\multicolumn{4}{p{11cm}}{\footnotesize $^\ddagger$\,Three equations
  (Coulomb's Law, Lorentz Force, Diffraction) have raw \rmse{}$=$\texttt{nan}
  in the source file despite good $\Rsq$ on the same equations, and one
  (Bernoulli's Equation) reports a $10^{150}$ sentinel/overflow value; all
  four are excluded from this mean. Over the 26 remaining equations, mean
  \rmse{}$\,\approx 2{,}740$. This \rmse{}-specific irregularity is
  unexplained and is noted for the record rather than resolved here.}\\
\end{tabular}
\end{table}

%% ============================================================
\section{Sweep Results: Noise Robustness (M3 vs.\ M4)}
\label{sec:noise}

\begin{table}[H]
\centering
\caption{$\Rsq$ statistics per noise level ($n=200$, 30 equations).
}
\label{tab:r2_noise}
\renewcommand{\arraystretch}{1.2}
\small
\begin{tabular}{l r r r r r r}
\toprule
& \multicolumn{3}{c}{\textcolor{m3blue}{\textbf{\EHD{} (M3)}}}
& \multicolumn{3}{c}{\textcolor{m4red}{\textbf{\HSL{} (M4)}}}\\
\cmidrule(lr){2-4}\cmidrule(lr){5-7}
$\sigma$ & Median & Min & Std & Median & Min & Std\\
\midrule
0\%    & 1.000000 & 0.999979 & $4.8\times10^{-6}$ & 1.0000000 & 1.000000 & $0.0$\\
0.05\% & 0.9999998 & 0.997348 & 0.0010 & 0.9999998 & 0.997807 & 0.0006\\
0.1\%  & 0.9999998 & 0.992748 & 0.0016 & 0.9999998 & 0.997810 & 0.0006\\
0.5\%  & 0.9999998 & 0.997504 & 0.0010 & 0.9999998 & 0.997789 & 0.0006\\
1\%    & 0.9999998 & 0.997325 & 0.0010 & 0.9999998 & 0.997720 & 0.0007\\
\bottomrule
\end{tabular}
\end{table}

\paragraph{Key observations.} The tables in this section (\ref{tab:r2_noise}, \ref{tab:rr_noise}, \ref{tab:time_noise}) are computed from the raw noise-sweep source files (\texttt{noise\_sweep\_20260812\_155519\_nshards03.json}, \texttt{...\_155520\_nshards05.json}, \texttt{...\_155532\_nshards01.json}, \texttt{...\_155541\_nshards02.json}, \texttt{...\_155544\_nshards04.json}), cross-checked against their raw per-test companions \texttt{protocol\_core\_noiseless\_20260812\_170344.json} and four \texttt{protocol\_core\_noisy\_*} files with matching timestamps to the second. The noise levels recorded in these files are $\sigma\in\{0, 0.05, 0.1, 0.5, 1\}\%$. M3 exhibits a std of $\sim$0.001--0.0016 across $\sigma>0$, higher
than M4's $\sim$0.0006--0.0007. M4 shows zero-variance convergence at
$\sigma=0\%$: all 30 equations reach $\Rsq=1.000000$ exactly, because the
LLM-first path generates the exact formula and the NN residual correction
drives the fit to machine precision. \textbf{M3 does not show the same
noiseless convergence}: at $\sigma=0\%$ M3's strict recovery rate
(threshold $\Rsq\geq0.999999$, per the source file's own recorded
\texttt{threshold\_used}) is \textbf{76.7\% (23/30)}, not 100\% --- see \S\ref{tab:rr_noise} below.

\begin{table}[H]
\centering
\caption{Recovery rate and catastrophic failure count per noise level ($n=200$).
Recovery uses each shard's recorded threshold: $\Rsq\geq0.999999$ at
$\sigma=0\%$, $\Rsq\geq0.995$ at $\sigma=0.05$, $0.1$ and $0.5\%$, and
$\Rsq\geq0.99$ at $\sigma=1\%$, so rates are not on a single scale across
rows.
}
\label{tab:rr_noise}
\small
\begin{tabular}{lrrrr}
\toprule
$\sigma$ & M3 Recovery & M3 Catastrophic & M4 Recovery & M4 Catastrophic\\
\midrule
0\%    & \textbf{76.7\%$^\ddagger$} & 0 & 100.0\% & 0\\
0.05\% & 100.0\%          & 0 & 100.0\% & 0\\
0.1\%  & 96.7\%           & 0 & 100.0\% & 0\\
0.5\%  & 100.0\%          & 0 & 100.0\% & 0\\
1\%    & 100.0\%          & 0 & 100.0\% & 0\\
\bottomrule
\multicolumn{5}{l}{\footnotesize$^\ddagger$ Source file \texttt{noise\_sweep\_..\_nshards01.json}, $\sigma=0$,}\\
\multicolumn{5}{l}{\footnotesize \texttt{threshold\_used}=0.999999: 23/30. M3 is \emph{not} perfect on noiseless recovery.}\\
\end{tabular}
\end{table}

\begin{table}[H]
\centering
\caption{Average per-equation computation time (seconds) per noise level.
Every cell is the per-equation \texttt{time} field averaged across all 30 tests in the five matching raw runner logs: \texttt{protocol\_core\_noiseless\_20260812\_170344.json} ($\sigma=0\%$) and \texttt{protocol\_core\_noisy\_20260812\_173425.json} ($\sigma=0.05\%$), \texttt{...\_172958.json} ($\sigma=0.1\%$), \texttt{...\_173523.json} ($\sigma=0.5\%$), and \texttt{...\_175318.json} ($\sigma=1\%$). The maximum M3 value across all 150 equation-runs (5 files $\times$ 30 equations) is 7.77\,s.}
\label{tab:time_noise}
\small
\begin{tabular}{lrrl}
\toprule
$\sigma$ & \textcolor{m3blue}{M3 avg (s)} & \textcolor{m4red}{M4 avg (s)} & Speedup\\
\midrule
0\%    & 7.55 & 6.92 & $1.1\times$\\
0.05\% & 7.66 & 6.55 & $1.2\times$\\
0.1\%  & 7.44 & 6.46 & $1.2\times$\\
0.5\%  & 7.77 & 6.56 & $1.2\times$\\
1\%    & 7.01 & 6.33 & $1.1\times$\\
\bottomrule
\end{tabular}
\end{table}

\paragraph{Note on M3 noiseless cost.} The source file for $\sigma=0\%$ records a mean M3 time of 7.55\,s, in line with every other noise level in \Cref{tab:time_noise}; none shows an exhaustive-search spike. The per-equation averages were re-derived from the raw runner logs (see the note on \Cref{tab:time_noise}) and agree with the table to two decimal places at all five noise levels. Across all
five real noise levels, M3 and M4 have comparable, stable per-equation
runtimes (M3: 7.0--7.8\,s; M4: 6.3--6.9\,s), with M4 consistently
$\sim$1.1--1.2$\times$ faster than M3.

%% ============================================================
\section{Sweep Results: Sample Complexity (M3 vs.\ M4)}
\label{sec:sc}

\begin{table}[H]
\centering
\caption{$\Rsq$ and RMSE per sample size (noiseless, 30 equations;
recovery thresholds used by the runner: $0.999$ at $n=50$, $0.9995$ at
$n=100$, $0.9999$ at $n\geq200$).
All six sample-size shards (\texttt{sample\_complexity\_20260812\_*\_nshards01--06.json}, $n=50$, $100$, $200$, $500$, $750$, $1000$, all \texttt{mode}: \texttt{noiseless}) were checked cell by cell against each shard's \texttt{per\_equation} data.}
\label{tab:sc_metrics}
\renewcommand{\arraystretch}{1.2}
\small
\begin{tabular}{r r r r r r r}
\toprule
& \multicolumn{3}{c}{\textcolor{m3blue}{\textbf{\EHD{} (M3)}}}
& \multicolumn{3}{c}{\textcolor{m4red}{\textbf{\HSL{} (M4)}}}\\
\cmidrule(lr){2-4}\cmidrule(lr){5-7}
$n$ & Med $\Rsq$ & Min $\Rsq$ & Med RMSE & Med $\Rsq$ & Min $\Rsq$ & Med RMSE\\
\midrule
  50 & 1.000000 & 0.9982 & 0.0006 & 1.0000000 & 1.0000 & 0.0000\\
 100 & 1.000000 & 0.9997 & 0.0011 & 1.0000000 & 1.0000 & 0.0000\\
 200 & 1.000000 & 1.0000 & 0.0009 & 1.0000000 & 1.0000 & 0.0000\\
 500 & 1.000000 & 0.9999 & 0.0016 & 1.0000000 & 1.0000 & 0.0000\\
 750 & 1.000000 & 0.9997 & 0.0021 & 1.0000000 & 1.0000 & 0.0000\\
1000 & 1.000000 & 0.9999 & 0.0013 & 1.0000000 & 1.0000 & 0.0000\\
\bottomrule
\end{tabular}
\end{table}

\HSL{} is exact (min $\Rsq=1.0000$, RMSE$=0.0000$) at every tested
sample size. \EHD{} reaches median $\Rsq=1.000000$ by $n=50$, but its
\emph{recovery rate} (fraction $\geq0.9999$) is 96.7\% (29/30) at
$n=50$, reaches 100\% at $n=100$ and $n=200$, then dips back to 96.7\%
at $n=500$, $750$, and $1000$ --- i.e.\ one equation stays just below
threshold at large $n$ rather than the metric monotonically improving
with sample size. ``Both converge at $n=50$'' is accurate for
\emph{median} $\Rsq$ but overstates \EHD's recovery-rate behaviour.

%% ============================================================
\section{Convergence Analysis}
\label{sec:convergence}

M3 shows stable convergence: $\Rsq$ std $\leq 0.001$ for $\sigma>0$.
M4 exhibits zero-variance convergence (std($\Rsq$) = 0.000) at all noise
levels in v2.

Both methods reach median $\Rsq=1.000000$ by $n=50$ in the noiseless sweep (all six sample-size
shards now verified; \Cref{tab:sc_metrics}).

\paragraph{Convergence boundary.}
The following conditions are needed for comprehensive generalisation claims:
(i)~higher-dimensional inputs ($\geq 6$ variables);
(ii)~PDE-form targets;
(iii)~extension to the full AI~Feynman dataset ($>100$ equations).

%% ============================================================
\section{Win Rate and Architecture Analysis}
\label{sec:winrate}

\begin{table}[H]
\centering
\caption{Head-to-head win rates (M3 vs.\ M4) on $\Rsq$. Noise sweep: 150 paired
comparisons (5 noise levels $\times$ 30 equations; pairs matched by shard and
test description). Sample-complexity (SC) sweep: 180 paired comparisons (6 sample
sizes $\times$ 30 equations; pairs matched by sample size and equation, from the
six \texttt{sample\_complexity\_*.csv} shards). A pair is \emph{tied} if both systems have
$\Rsq>0.9999$ (the paper's original tie rule); otherwise the higher $\Rsq$ wins.
}
\label{tab:winrate}
\small
\begin{tabular}{l r r}
\toprule
\textbf{Outcome} & \textbf{Noise sweep} & \textbf{SC sweep}\\
\midrule
M3 strictly higher $\Rsq$ & 0/150 (0.0\%) & 0/180 (0.0\%)\\
M4 strictly higher $\Rsq$ & 28/150 (18.7\%) & 8/180 (4.4\%)\\
Tied (both $\Rsq>0.9999$)  & 122/150 (81.3\%) & 172/180 (95.6\%)\\
\bottomrule
\end{tabular}
\end{table}

\paragraph{Sensitivity to the tie rule.}
The split depends strongly on how ties are defined (noise sweep, $n=150$).
Under exact floating-point comparison M3 is nominally higher on 88 pairs and M4 on
62, but most of these differences are below $10^{-9}$ and reflect numerical noise
rather than a real accuracy gap. With an absolute tolerance of $10^{-6}$, M3 is
higher on 0, M4 on 35 (23.3\%) and 115 are tied (76.7\%); at $10^{-4}$ the figures
are 0 / 23 (15.3\%) / 127 (84.7\%); at $10^{-3}$ they are 0 / 19 (12.7\%) /
131 (87.3\%). Under no tolerance of $10^{-6}$ or looser does M3 exceed M4 on any
pair, and M4 is strictly higher on 12.7--23.3\% of pairs depending on the
tolerance. M4 therefore never loses to M3 beyond numerical noise on the noise
sweep, and most comparisons are ties.

The sample-complexity sweep ($n=180$) shows the same pattern. Under exact
comparison M3 is higher on 0 pairs, M4 on 179 and 1 is tied; with an absolute
tolerance of $10^{-6}$ the figures are 0 / 38 (21.1\%) / 142 (78.9\%); at
$10^{-4}$ they are 0 / 8 (4.4\%) / 172 (95.6\%); at $10^{-3}$ they are 0 / 1
(0.6\%) / 179 (99.4\%). M4's 8 strict wins under the $\Rsq>0.9999$ rule fall at
$n=50$ (2), $n=100$ (3), $n=200$ (0), $n=500$ (1), $n=750$ (1) and $n=1000$ (1).
M3 is never higher than M4 under any tie rule, including exact comparison.

\begin{table}[H]
\centering
\caption{Internal routing statistics --- evidence that M3 and M4 are genuinely
         distinct architectures. From the five \texttt{protocol\_core} result
         files (one noiseless and four noisy runs,
         $\sigma\in\{0,0.05,0.1,0.5,1\}\%$; $5$ noise levels $\times$ $30$
         equations $= 150$ records per system). ``Ensemble-first'' (M3) and
         ``LLM-first'' (M4) denote the share of records whose recorded
         \texttt{metadata.decision} routed to that branch; this is a routing
         share, not a solve rate, and the \texttt{success} flag is not used.
         The pool mixes one noiseless run ($\sigma=0$) with four noisy runs.}
\label{tab:arch}
\small
\begin{tabular}{L{3.5cm} r r}
\toprule
\textbf{Decision path} & \textbf{M3 (\EHD)} & \textbf{M4 (\HSL)}\\
\midrule
Primary path (routing share)
                & 140/150 (93.3\%) ensemble-first & 123/150 (82.0\%) LLM-first\\
Alternative path
                & 10/150 (6.7\%) NN              & 27/150 (18.0\%) ensemble\\
Secondary path  & NN refinement                  & NN residual correction\\
Architecture type & Ensemble-first               & LLM-first + NN-refinement\\
\bottomrule
\end{tabular}
\end{table}

\begin{table}[H]
\centering
\caption{Routing by noise level (30 equations per level), same source as
         \Cref{tab:arch}. M3 routes identically at every level; M4 routes all
         noiseless records to the LLM but $20$--$27\%$ of noisy records to the
         ensemble.}
\label{tab:arch_by_noise}
\small
\begin{tabular}{l r r r r}
\toprule
$\sigma$ (\%) & M3 ensemble & M3 NN & M4 LLM & M4 ensemble\\
\midrule
0      & 28 & 2 & 30 & 0\\
0.05   & 28 & 2 & 22 & 8\\
0.1    & 28 & 2 & 24 & 6\\
0.5    & 28 & 2 & 23 & 7\\
1      & 28 & 2 & 24 & 6\\
\midrule
Total  & 140 & 10 & 123 & 27\\
\bottomrule
\end{tabular}
\end{table}

\paragraph{When to prefer M4.}
\begin{itemize}
\item \textbf{Prefer M4} when wall-clock time is a factor
  (M4 is $\sim$1.1--1.2$\times$ faster than M3; \S\ref{sec:noise}), near-zero RMSE is prioritised, or $\sigma>0$ conditions are
  expected.
\end{itemize}

%% ============================================================
\section{Root-Cause Analysis: Newton's Gravity Measurement Bug}
\label{sec:bugfix}

\subsection{Symptom}

M4 reported $\Rsq \in [0.42, 0.87]$ on Newton's gravitational force across
all 11 tested conditions in v1, despite the LLM generating the correct formula
on every run.

\subsection{Root Cause}

\texttt{evaluate\_llm\_formula} used:
\begin{verbatim}
r2 = 1 - (ss_res / ss_tot) if ss_tot > 1e-10 else 0
\end{verbatim}
For Newton's gravity at tested parameter ranges, outputs are $\sim 10^{-11}$\,N,
giving $\mathrm{ss\_tot} \approx 1.8 \times 10^{-18}$ — eight orders of
magnitude below the threshold.  The function always returned
\texttt{r2 = 0.0}, triggering NN fallback.

\subsection{Fix Applied}

\begin{verbatim}
_tol = 1e-10 * float(np.max(np.abs(y_true))**2) * len(y_true)
r2   = 1 - (ss_res / ss_tot)
        if ss_tot > max(_tol, 1e-300) else 0
\end{verbatim}

The same fix was applied at three code locations:
\texttt{evaluate\_llm\_formula} (line 462),
\texttt{train\_nn} (line 384), and
the ensemble blend $\Rsq$ calculation (line 700).

\subsection{Additional Fixes in the Same Commit}

\begin{itemize}
  \item \texttt{force\_llm} flag: now correctly bypasses NN training.
  \item \texttt{PureLLMBaseline}: validates non-empty \texttt{python\_code}.
  \item Training epochs: corrected from 300 to 1000 (scheduler requires $\geq 1000$).
  \item Log-space guard: unconditionally enables log-space training when
    \texttt{y\_range > 3.0} decades.
\end{itemize}

\paragraph{General contribution.}
Any SR benchmark harness using absolute sum-of-squares thresholds will
silently misclassify correct formulas on equations with output values far
from unity.  The relative threshold fix should be adopted as standard practice.

%% ============================================================
\section{Discussion}
\label{suppb:discussion}

\subsection{Hardcoding and LLM Benchmark Integrity}

Of \PureLLM{}'s 30 noiseless rows, 11 (36.7\%) did not call the language model: 8 return the benchmark's own
ground-truth expression and 3 use a fixed form with parameters fitted to the data
(\Cref{tab:hardcoded}). In the 170344 file, \EHD{}'s returned formula is character-for-character identical to \PureLLM{}'s on all 11 of these bypassed rows, and on only 2 of the 19 live rows. The cause is in the code: \EHD{} delegates formula generation to \texttt{PureLLMBaseline.generate\_formula}, which tries the hardcoded formulas before any language-model request; \EHD{}'s pass count on the bypassed rows (6 of 11, against 17 of 19 on live rows) also depends on its NN and ensemble logic. The all-domains LLM+NN hybrid delegates in the same way, but its stored output has no formula, so whether the bypass fired for it cannot be checked from the JSON. \PureLLM{} passes 29 of 30 rows (96.7\%) overall, 11 of 11 bypassed rows and 18 of 19
(94.7\%) \emph{live} rows, i.e.\ rows where the model was actually called. The live figure is not an unbiased
estimate of unaided performance: according to comments in the baseline source, the bypass branches were added for
equations on which the model had repeatedly failed, so excluding them favours \PureLLM{}. Counting every bypassed
row as a failure gives 18/30, so the unaided pass count lies between 18/30 and 29/30; the baseline has not been rerun
with the bypass disabled.
The \PureLLM{} figure in \Cref{tab:overall} (29/30; 18/19 on live rows) is the recount over the per-domain files; the noise tables (\Cref{app:noise_sensitivity}) use the \texttt{170344} file, which reports 30/30. The one-test difference is Curie's law, which fails in the per-domain recount and passes in the \texttt{170344} run.
Explicit hardcoding audits should be standard in any
benchmark involving LLM-based SR.

\subsection{M3 vs.\ M4 Trade-off}

With the measurement bug resolved, M3 and M4 are \textbf{close to
equivalent on $\Rsq$}: both reach median $\Rsq = 1.000000$ at every
tested noisy condition and sample size, though M3's \emph{recovery
rate} dips to 96.7\% at $\sigma=0.1\%$ (\Cref{tab:rr_noise}) and at
$n=500/750/1000$ (\Cref{tab:sc_metrics}), while M4 is 100\% throughout.
The distinction is that M4 is efficiency-optimised ($\sim$1.1--1.2$\times$ faster,
near-zero RMSE).

Under the $\Rsq>0.9999$ rule the tie rate is $81.3\%$ on the noise
sweep (122/150) and $95.6\%$ on the sample-complexity sweep (172/180), or
$89.1\%$ over the 330 pairs combined (\Cref{tab:winrate}). Direct computation
confirms that, on the noise sweep, M4 never loses to M3 on $\Rsq$ beyond
numerical noise ($|\Delta\Rsq|\le10^{-6}$), and that on the sample-complexity
sweep M3 is never higher than M4 under any tie rule.

%% ============================================================
\section{Limitations}
\label{suppb:limitations}

\begin{enumerate}
  \item \textbf{Benchmark scope.} 30 equations, 2--5 input variables.
    Higher-dimensional and PDE targets remain untested.
  \item \textbf{Memorisation scope.} The hardcoding audit detects verbatim
    recall but not partial memorisation.
  \item \textbf{Baseline breadth.} Only M3 and M4 covered by the sweep;
    adding PySR and DSR sweeps would strengthen generalisation claims.
  \item \textbf{M3 noiseless recovery.}  At $\sigma=0\%$, M3's strict recovery rate is 76.7\%
    (23/30, \Cref{tab:rr_noise}), not 100\%.
  \item \textbf{No formula complexity metric.} Expression tree depth and
    operator count should accompany $\Rsq$ in future evaluations.
\end{enumerate}

%% ============================================================
\section{Conclusions}
\label{sec:conclusions}

\begin{enumerate}
  \item \textbf{Both M3 and M4 achieve 100\% recovery} at all tested noisy
    conditions ($0.05$--$1\%$ noise) and all
    sample sizes ($50$--$1000$).
  \item \textbf{M4 reaches 100\% at $\sigma=0$} (noiseless, v2-corrected);
    \textbf{M3 reaches only 76.7\% (23/30) at $\sigma=0$}
    (\Cref{tab:rr_noise}). \HDS{} (M6, tools track) trails at 70.0\% (21/30)
    under the separate six-method noiseless protocol (30 equations;
    \Cref{tab:overall}), with nine failing equations (Arrhenius rate
    constant, Nernst equation, dielectric polarisation, capacitor energy,
    Coulomb's law, kinetic energy, Gaussian/normal PDF, Bose--Einstein
    occupation number, Fermi--Dirac occupation number).

  \item \textbf{M4 is $\sim$1.1--1.2$\times$ faster} with near-zero RMSE
    ; the throughput advantage is modest, not
    decisive.
  \item \textbf{The Newton failure in v1 was a measurement bug}: absolute
    $\Rsq$ threshold misclassified correct formulas for tiny-scale equations.
    The relative threshold fix is a general contribution to SR benchmarking.
  \item \textbf{Three methodological issues corrected}: undisclosed hardcoding,
    scale-biased \rmse, mismatched thresholds.
\end{enumerate}

%% ── Bibliography ─────────────────────────────────────────────────────────────
\bibliography{references}

%% ============================================================
\appendix

\section{Reproducibility Figure Generation Inventory}
\label{app:figures}

\noindent This appendix documents the auxiliary figures produced by
\texttt{generate\_figures.py} for reproducibility purposes. It is
distinct from the figures actually embedded in the compiled
manuscript (5 files).

\begin{table}[H]
\centering
\caption{Complete figure inventory; all files in \texttt{figures/}.}
\label{tab:figlist}
\small
\renewcommand{\arraystretch}{1.25}
\begin{tabular}{L{4.8cm} L{1.5cm} L{1.8cm} L{3.2cm}}
\toprule
\textbf{Filename} & \textbf{Sweep} & \textbf{Section} & \textbf{Content}\\
\midrule
\texttt{fig1\_r2\_vs\_noise.pdf}       & Noise & \ref{sec:noise}        & Median $\Rsq$ vs $\sigma$\\
\texttt{fig2\_rmse\_vs\_noise.pdf}     & Noise & \ref{sec:noise}        & Median RMSE vs $\sigma$\\
\texttt{fig3\_time\_vs\_noise.pdf}     & Noise & \ref{sec:noise}        & Avg time vs $\sigma$\\
\texttt{fig4\_r2\_vs\_n.pdf}           & SC    & \ref{sec:sc}           & Median $\Rsq$ vs $n$\\
\texttt{fig5\_rmse\_vs\_n.pdf}         & SC    & \ref{sec:sc}           & Median RMSE vs $n$\\
\texttt{fig6\_time\_vs\_n.pdf}         & SC    & \ref{sec:sc}           & Median time vs $n$\\
\texttt{fig7\_recovery\_vs\_noise.pdf} & Noise & \ref{sec:noise}        & Recovery rate vs $\sigma$\\
\texttt{fig8\_recovery\_vs\_n.pdf}     & SC    & \ref{sec:sc}           & Recovery rate vs $n$\\
\texttt{fig9\_minr2\_vs\_noise.pdf}    & Noise & \ref{sec:noise}        & Min $\Rsq$ vs $\sigma$\\
\texttt{fig10\_r2\_boxplot\_noise.pdf} & Noise & \ref{sec:convergence}  & Per-eq $\Rsq$ box plots\\
\texttt{fig11\_recovery\_heatmap.pdf}  & Both  & \ref{sec:convergence}  & Recovery heatmap $\sigma \times n$\\
\midrule
\texttt{fig\_runtime\_comparison.png}  & ---   & \ref{sec:noiseless}    & Runtime bar chart, 6 methods\\
\texttt{fig\_comparative\_table.png}   & ---   & \ref{sec:noiseless}    & Domain $\times$ method table\\
\bottomrule
\end{tabular}
\end{table}

\noindent\textbf{Relationship to the submitted PDF:} the 13 files
above are reproducibility documentation only and are not
\verb|\includegraphics|'d anywhere in this document or in the main
manuscript. The compiled manuscript embeds a separate, disjoint set
of 5 figures directly via \verb|\includegraphics| (see
\texttt{jmlr\_paper\_main.tex}). Only those 5 files are bundled in
\texttt{submission/sources/figures/}; the 13 listed here should be
regenerated from source by anyone who wants them, not shipped in the
submission package.

\section{Method Abbreviations}
\label{app:abbrev}

\begin{tabular}{L{1.8cm} L{2.5cm} L{6.0cm}}
\toprule
\textbf{Short} & \textbf{Code} & \textbf{Full name}\\
\midrule
\PureLLM & M1 & PureLLM Baseline (core)\\
\INN     & M2 & ImprovedNN (core)\\
\EHD     & M3 & EnhancedHybridSystemDeFi (core)\\
\HSL     & M4 & HybridSystemLLMNN all-domains (core)\\
\SEL     & M5 & SymbolicEngineWithLLM (tools)\\
\HDS     & M6 & HybridDiscoverySystem v50\_2 (tools)\\
\bottomrule
\end{tabular}

\section{PureLLM Required Disclosure}
\label{app:disclosure}

Required in any JMLR submission including \PureLLM{} results:

\begin{quote}
\textit{The PureLLM Baseline does not call the language model for 11 of 30
benchmark equations (\texttt{is\_hardcoded: true}): for 8 it returns the benchmark's own
ground-truth formula and for 3 a fixed form whose parameters are fitted to the data by
least squares. Its pass rate is 29/30 over all rows and 18/19 on the 19 live rows
(where the model was called). The live figure favours PureLLM because the bypass
was added for equations on which the model had failed, so the unaided pass count
lies between 18/30 and 29/30; the baseline has not been rerun with the bypass disabled.}

\end{quote}

% ====================================================================
% NEW APPENDICES (extracted from jmlr_paper_main.tex)
% ====================================================================
\section{Benchmark Reproducibility Statement}
\label{app:repro_statement}

\subsection*{Code, Data, and Protocol}

All experimental results in the main paper's \emph{Benchmark Design} section were produced
using \texttt{experiment\_protocol\_benchmark\_v2.py} (version~2.0,
2026-03-02), which will be released alongside the final paper.
The protocol file encodes all data-generation functions, random seeds, noise
levels, and recovery thresholds as class constants and is self-documenting:

\begin{verbatim}
python experiment_protocol_benchmark_v2.py --describe
\end{verbatim}

\noindent The 30-equation benchmark is defined in
\texttt{\_build\_feynman\_equations()} (lines 480--953 of the protocol file; the
function name uses the label ``feynman'', but the set is custom, not a Feynman subset).
Each equation specifies its identifier, ground-truth formula, variable
ranges, and generating function; data are produced analytically, eliminating
external dataset dependencies for the physics equations.

\subsection*{Exact Reproduction Commands}

\paragraph{Phase~2 (noisy, $R^2>0.995$).}
\begin{verbatim}
python run_comparative_suite_benchmark_v2.py \
    --methods 1 --samples 200 --no-llm-cache
\end{verbatim}
Output: \texttt{protocol\_core\_noisy\_<timestamp>.json}.
The \texttt{--no-llm-cache} flag forces independent fresh API calls for all
30 Pure~LLM evaluations; this was verified via cache-clearing.

\paragraph{Phase~3 (noiseless, $R^2\geq0.999999$, Table~\ref{tab:overall}).}
The command and output below come from the one verified noiseless source file, \texttt{protocol\_core\_noiseless\_20260812\_170344.json}, whose internal \texttt{script} field is quoted directly.
\begin{verbatim}
python run_protocol_benchmark_core.py \
    --noiseless --threshold 0.999999     \
    --nn-seeds 3 --samples 200           \
    --method-timeout 900                 \
    --pysr-timeout 1100
\end{verbatim}
Script version: \texttt{run\_protocol\_benchmark\_core.py v2.2 (sign-fix +
log-widen + domain-guard + formula-hash + complexity-score + json-export +
extrap\_r2\_far-fix)}, per the source file's own \texttt{script} field.\\
Output: \texttt{protocol\_core\_noiseless\_20260812\_170344.json}
(timestamp 2026-08-12T17:03:44; 30 equations across 11 non-Feynman domains
--- biology, chemistry, electrochemistry, electromagnetism, electrostatics,
magnetism, mechanics, optics, probability, quantum, thermodynamics --- run
against all 6 methods, with per-method pass counts as read from that file, which
\Cref{tab:overall} reproduces except for the \PureLLM{} pass count (29/30 in the table, from the per-domain recount; 30/30 in this file); the other counts in this file are \SEL{} 28/30, \HSL{} 30/30, \EHD{} 23/30, \HDS{} 21/30, \INN{} 0/30; wall-clock and hardware not recorded in
the source file, so no timing claim is made here.

\textbf{Note on timeout upgrade:} \texttt{--method-timeout} was set to 300\,s
for tests 1--18 and upgraded to 900\,s for tests 19--30. Arrhenius (test~4)
ran under the 300\,s limit and the harness recorded a 27574\,s wall-clock
before the run ended; this is recorded as a genuine failure independent of
the v2 bug fix. PySR is
invoked via \texttt{juliacall}, which embeds Julia in-process rather than
spawning a separate OS process, so there was no external Julia PID being
``killed.'' The 27574\,s figure instead reflects a gap in the harness's own
timeout wiring: the outer \texttt{ThreadPoolExecutor} timeout fires at
300\,s but abandons the background thread without a reference back to the
underlying call, which then continues to run unobserved.

Nernst (test~6) also timed out under the 300\,s default; counted as failure.

\subsection*{Random Seeds}

\begin{itemize}
  \item Base seed: \texttt{42} (class constant
    \texttt{BenchmarkProtocol.seed = 42}, line~1091 of the protocol file).
    Matches the Core~15 benchmark seed for cross-campaign consistency.
  \item Per-equation seed: \texttt{seed + offset $\times$ 7}, where
    \texttt{offset} is the 0-based index in the equation list
    (line~1208: \texttt{eq.generate(n, noise, seed=s + offset * 7)}).
    This ensures each of the 30 equations receives a distinct, reproducible
    seed regardless of call order.
  \item Neural network seeds (Phase~3): 3 independent seeds per equation;
    median $R^2$ reported. Exact seed values are logged in the output JSON.
\end{itemize}

\subsection*{Software Environment}

\begin{table}[H]
\centering
\caption{Software and hardware environment for benchmark runs
         (verified 18 March 2026 via \texttt{hardware\_info.py} and
         \texttt{pip show}).}
\label{tab:software_env}
\small
\begin{tabular}{ll}
\toprule
\textbf{Component} & \textbf{Version / details} \\
\midrule
Python           & 3.12.2 \\
PySR             & 1.5.9 \\
Julia runtime    & 1.12.2 \\
juliacall        & 0.9.26 (Python$\leftrightarrow$Julia bridge used by PySR) \\
NumPy            & 2.3.5 \\
SciPy            & 1.16.3 \\
SymPy            & 1.14.0 \\
Pandas           & 2.3.3 \\
Matplotlib       & 3.10.7 \\
Seaborn          & 0.13.2 \\
scikit-learn     & 1.7.2 \\
Anthropic SDK    & 0.73.0 (\texttt{claude-sonnet-5},
                   temperature~0.25) \\
\midrule
Hardware         & Intel Celeron T3100 @ 1.90\,GHz (2~cores);
                   3\,GB DDR2; no GPU \\
OS               & Kali GNU/Linux Rolling \\
PySR populations & 2 (hardware limit: 2-core CPU) \\
\bottomrule
\end{tabular}
\end{table}

\noindent\textbf{Note on Anthropic SDK version compatibility.}
\texttt{anthropic==0.73.0} is substantially newer than early-2024 releases.
The Messages API interface is stable across 0.18$+$ versions; earlier
installations may require minor call-site adjustments but the prompts and
response parsing are otherwise unchanged.

\noindent PySR uses Julia under the hood; the 1100\,s
\texttt{--pysr-timeout} (\texttt{FEYNMAN\_TIMEOUT}, the inner PySR fit
budget---distinct from the 900\,s outer \texttt{--method-timeout} above,
which governs the per-method wall-clock allowance, not the PySR fit itself)
was required to accommodate Julia JIT initialisation on the first run.
Subsequent equations in the same process are faster. The 27574\,s Arrhenius
figure is a harness-timeout artifact, not evidence that PySR's complexity
search fails to converge on single-variable flat-landscape equations; see
the note above.

\subsection*{PureLLM Disclosure}

Following~\citet{memorization2024audit}, each PureLLM noiseless row was audited for whether the
language model was called. \textbf{11 of 30 PureLLM rows bypass the model
(\texttt{is\_hardcoded=True})}: Michaelis--Menten, Logistic growth,
Allometric scaling, Nernst, Clausius--Mossotti, Lorentz force, Gaussian PDF,
Zeeman energy, Bose--Einstein, Planck radiance, and Fourier heat. Eight of these
return the benchmark's ground-truth formula; Michaelis--Menten, Logistic growth and Allometric scaling use a fixed form with
least-squares parameters fitted to the data. PureLLM passes 29/30 over all rows and 18/19 on the 19 live rows
(where the model was called); because the bypass was added for equations on which the model had failed,
the live figure favours PureLLM and the unaided pass count lies between 18/30 and 29/30.

PureLLM is listed in all tables, and the main paper reports its comparisons with and without the bypassed rows.

A controlled re-run with \texttt{--no-llm-cache} confirmed that all
30 equations issued independent API calls and produced identical scores
(median $R^2=0.9976$, std $=2.5\times10^{-4}$), ruling out any
caching artefact.

\section{Hyperparameter Settings}
\label{app:hyperparameters}

\begin{table}[H]
\centering
\caption{PySR Configuration for Symbolic Regression}
\label{tab:hyperparameters}
\begin{tabular}{@{}lc@{}}
\toprule
\textbf{Parameter} & \textbf{Value} \\
\midrule
Population size & 100 \\
Generations & 500 \\
Populations (parallel) & 2 (hardware limit: 2-core CPU) \\
Crossover probability & 0.8 \\
Mutation probability & 0.1 \\
Complexity penalty & 0.01 \\
Tournament size & 3 \\
Elitism & 5 \\
Max expression depth & 15 \\
Optimization iterations & 50 \\
\midrule
\textbf{Validation Thresholds} & \\
\midrule
Acceptance $R^2$ (Criterion 1) & 0.99 \\
Validation score (Criterion 1) & 30 \\
Acceptance $R^2$ (Criterion 2) & 0.95 \\
Validation score (Criterion 2) & 80 \\
Extrapolation $R^2$ & 0.80 \\
Numerical stability ($\alpha$) & 0.95 \\
\bottomrule
\end{tabular}
\end{table}

\begin{table}[H]
\centering
\caption{Complete Hyperparameter Specifications for All Methods}
\label{tab:hyperparameters_complete}
\small
\begin{tabular}{lll}
\toprule
\textbf{Method} & \textbf{Parameter} & \textbf{Value} \\
\midrule
\multirow{8}{*}{Neural Network}
& Architecture & [64, 32, 1] \\
& Activation & ReLU \\
& Optimizer & Adam \\
& Learning rate & 0.01 \\
& Batch size & Full batch (160 samples) \\
& Epochs & 200 \\
& Weight initialization & Xavier uniform \\
& Dropout & None \\
\midrule
\multirow{7}{*}{Pure LLM}
& Model & claude-sonnet-5 \\
& Temperature & 0.25 \\
& Max tokens & 1024 \\
& Top-p & 1.0 \\
& Few-shot examples & 2-3 per domain \\
& Prompt template & Domain-specific \\
& Retry on failure & 3 attempts \\
\midrule
\multirow{12}{*}{Hybrid v50\_2}
& SR iterations & 50 \\
& SR populations & 2 (hardware limit: 2-core CPU) \\
& Binary operators & [+, -, *, /] \\
& Unary operators & [exp, log, sqrt, sin, cos] \\
& Parsimony coefficient & 0.01 \\
& Tournament size & 10 \\
& Mutation rate & 0.1 \\
& Crossover rate & 0.9 \\
& LLM candidates & 5 \\
& LLM temperature & 0.25 \\
& Max complexity & 30 \\
& Validation threshold ($R^2$) & 0.99 \\
\bottomrule
\end{tabular}
\end{table}

% ============================================================================
% APPENDIX C: DETAILED EXPERIMENTAL RESULTS
% ============================================================================


\section{Detailed Experimental Results}
\label{app:detailed_results}

\subsection{Pure LLM Baseline: Per-Case Results}

% Auto-generated by tables/generate_tables.py
% Source: hypatiax/data/results/comparison_results/noise-noiseless/noiseless/defi/hypatiax_defi_benchmark_v3_results_seed42.json
% Date:   2026-10-08 10:53
% DO NOT EDIT MANUALLY — re-run 'make tables' to regenerate


\begingroup\small
\begin{longtable}{L{5.4cm}rr}
\caption{Pure LLM performance by test case (pinned DeFi benchmark, seed 42). \texttt{---}: the result file holds no numeric $R^2$ for that case.}\label{tab:llm_detailed}\\
\toprule
\textbf{Case} & \textbf{$R^2$} & \textbf{Time (s)} \\
\midrule
\endfirsthead
\toprule
\textbf{Case} & \textbf{$R^2$} & \textbf{Time (s)} \\
\midrule
\endhead
\bottomrule
\endfoot
Value at Risk at 95\% & 1.0000 & 8.23 \\
Value at Risk at 99\% & 1.0000 & 6.84 \\
Partial liquidation amount & 1.0000 & 12.01 \\
Constant product formula & 1.0000 & 6.71 \\
Simple Staking APY & 1.0000 & 5.95 \\
Loan-to-Value & $-1.41\times10^{5}$ & 7.61 \\
Spot price from AMM & 1.0000 & 6.79 \\
LP share percentage & $-1.8739$ & 5.94 \\
Long position unrealized PnL & 1.0000 & 4.97 \\
Short position unrealized PnL & 1.0000 & 4.74 \\
Funding rate cost & 1.0000 & 7.26 \\
Validator commission adjusted & 1.0000 & 5.69 \\
Slashing penalty & --- & 14.26 \\
Protocol reserve accumulation & 1.0000 & 7.09 \\
Leveraged position notional & 1.0000 & 6.84 \\
Cross-margin available balance & 1.0000 & 8.64 \\
Realized PnL for long & 1.0000 & 4.88 \\
LP fee earnings & 1.0000 & 8.81 \\
Multi-day Value at Risk & 1.0000 & 7.02 \\
ES scaling for multi-day & 1.0000 & 7.66 \\
Annualised Portfolio tracking error & 1.0000 & 7.26 \\
Incremental VaR & 1.0000 & 8.15 \\
Reserve ratio & 1.0000 & 6.63 \\
Impermanent loss breakeven fee rate & 1.0000 & 11.84 \\
Liquidation Price Long & --- & 7.56 \\
Liquidation Price Short & --- & 6.56 \\
Constant Product Price Impact & 1.0000 & 7.78 \\
Effective Leverage & $-12.4502$ & 14.12 \\
Borrowing Interest & 1.0000 & 6.67 \\
Compounding Staking Returns & 1.0000 & 7.21 \\
Portfolio Sharpe Ratio & 1.0000 & 5.50 \\
APY calculation & 1.0000 & 8.62 \\
Capital efficiency & --- & 8.40 \\
AMM output amount & 1.0000 & 10.10 \\
Price slippage percentage & $-3.2163$ & 6.97 \\
Information ratio & 1.0000 & 5.33 \\
Utilization rate of DeFi & --- & 6.64 \\
Borrow APY from utilization & 1.0000 & 6.45 \\
Supply APY from borrow & 1.0000 & 6.38 \\
Health factor & 1.0000 & 9.74 \\
Impermanent loss percentage & $-2.1798$ & 7.72 \\
Options Delta & --- & 16.66 \\
Collateral ratio & --- & 7.55 \\
Funding rate cost (extended) & 1.0000 & 7.88 \\
Staking reward for fixed lock-up & $-0.9246$ & 7.13 \\
Expected Shortfall at 95\% & 1.0000 & 7.17 \\
Expected Shortfall at 99\% & 1.0000 & 7.20 \\
Optimal LP Position (Kelly) & --- & 9.26 \\
Concentrated liquidity position width & $-2.59\times10^{8}$ & 11.11 \\
Portfolio VaR for two correlated & 1.0000 & 9.19 \\
Position margin ratio & 1.0000 & 6.19 \\
Concentrated liquidity position width (v2) & 1.0000 & 7.66 \\
Spot price from AMM reserve & 0.9810 & 10.03 \\
Black-Scholes Call Price & --- & 21.86 \\
Black-Scholes Put Price & --- & 16.19 \\
Component ES & $-0.2749$ & 9.61 \\
Gamma of option & --- & 12.63 \\
Vega of option & --- & 9.73 \\
Multi-Collateral LTV & 1.0000 & 8.55 \\
Correlated Portfolio VaR & 1.0000 & 8.41 \\
Impermanent loss in constant product & 1.0000 & 9.08 \\
Constant product formula (multivariate) & 1.0000 & 5.91 \\
Convexity Adjustment & 0.9998 & 18.49 \\
Liquidation price for leveraged long & --- & 10.06 \\
Liquidation price for leveraged short & --- & 6.68 \\
Maximum safe leverage & --- & 7.92 \\
Required collateral & --- & 8.04 \\
Forward price for derivative & 1.0000 & 6.74 \\
Portfolio Expected Shortfall for correlated & --- & 4.99 \\
Call option intrinsic & $-3.5798$ & 4.68 \\
Put-call parity & 1.0000 & 7.33 \\
Simple options moneyness & $-50.5674$ & 4.92 \\
Uniswap V3 virtual & $-9.41\times10^{6}$ & 14.65 \\
Theta of option & --- & 21.12 \\
\end{longtable}
\endgroup


\subsection{DeFi-Optimized HypatiaX: Per-Case Results}

% Auto-generated by tables/generate_tables.py
% Source: hypatiax/data/results/comparison_results/noise-noiseless/noiseless/defi/hypatiax_defi_benchmark_v3_results_seed42.json
% Date:   2026-10-08 10:53
% DO NOT EDIT MANUALLY — re-run 'make tables' to regenerate


\begingroup\small
\begin{longtable}{L{5.0cm}rrl}
\caption{HypatiaX performance by test case (pinned DeFi benchmark, seed 42). $R^2$ is that of the sub-method named by the routing decision, not hybrid's self-reported value; \texttt{---}: that sub-method returned no numeric $R^2$. Decisions without an independent baseline (\texttt{ensemble}, \texttt{v4\_residual\_nn}) show hybrid's own value.}\label{tab:defi_detailed}\\
\toprule
\textbf{Case} & \textbf{$R^2$} & \textbf{Time (s)} & \textbf{Decision} \\
\midrule
\endfirsthead
\toprule
\textbf{Case} & \textbf{$R^2$} & \textbf{Time (s)} & \textbf{Decision} \\
\midrule
\endhead
\bottomrule
\endfoot
Value at Risk at 95\% & 1.0000 & 1.34 & \texttt{llm} \\
Value at Risk at 99\% & 1.0000 & 1.08 & \texttt{llm} \\
Partial liquidation amount & 1.0000 & 1.13 & \texttt{llm} \\
Constant product formula & 1.0000 & 1.30 & \texttt{llm} \\
Simple Staking APY & 1.0000 & 1.30 & \texttt{llm} \\
Loan-to-Value & $-1.41\times10^{5}$ & 1.27 & \texttt{llm} \\
Spot price from AMM & 1.0000 & 1.07 & \texttt{llm} \\
LP share percentage & $-1.8739$ & 1.11 & \texttt{llm} \\
Long position unrealized PnL & 1.0000 & 1.26 & \texttt{llm} \\
Short position unrealized PnL & 1.0000 & 1.78 & \texttt{llm} \\
Funding rate cost & 1.0000 & 1.12 & \texttt{llm} \\
Validator commission adjusted & 1.0000 & 1.22 & \texttt{llm} \\
Slashing penalty & --- & 1.46 & \texttt{llm} \\
Protocol reserve accumulation & 1.0000 & 1.29 & \texttt{llm} \\
Leveraged position notional & 1.0000 & 1.05 & \texttt{llm} \\
Cross-margin available balance & 1.0000 & 1.20 & \texttt{llm} \\
Realized PnL for long & 1.0000 & 1.41 & \texttt{llm} \\
LP fee earnings & 1.0000 & 1.30 & \texttt{llm} \\
Multi-day Value at Risk & 1.0000 & 1.33 & \texttt{llm} \\
ES scaling for multi-day & 1.0000 & 1.25 & \texttt{llm} \\
Annualised Portfolio tracking error & 1.0000 & 1.21 & \texttt{llm} \\
Incremental VaR & 1.0000 & 2.09 & \texttt{llm} \\
Reserve ratio & 1.0000 & 1.35 & \texttt{llm} \\
Impermanent loss breakeven fee rate & 1.0000 & 1.59 & \texttt{llm} \\
Liquidation Price Long & --- & 1.18 & \texttt{llm} \\
Liquidation Price Short & --- & 1.28 & \texttt{llm} \\
Constant Product Price Impact & 1.0000 & 1.78 & \texttt{llm} \\
Effective Leverage & $-12.4502$ & 1.87 & \texttt{llm} \\
Borrowing Interest & 1.0000 & 1.26 & \texttt{llm} \\
Compounding Staking Returns & 1.0000 & 1.26 & \texttt{llm} \\
Portfolio Sharpe Ratio & 1.0000 & 1.13 & \texttt{llm} \\
APY calculation & 1.0000 & 1.10 & \texttt{llm} \\
Capital efficiency & --- & 1.68 & \texttt{llm} \\
AMM output amount & 1.0000 & 1.16 & \texttt{llm} \\
Price slippage percentage & $-3.2163$ & 1.27 & \texttt{llm} \\
Information ratio & 1.0000 & 1.18 & \texttt{llm} \\
Utilization rate of DeFi & --- & 1.69 & \texttt{llm} \\
Borrow APY from utilization & 1.0000 & 1.27 & \texttt{llm} \\
Supply APY from borrow & 1.0000 & 1.17 & \texttt{llm} \\
Health factor & 1.0000 & 1.57 & \texttt{llm} \\
Impermanent loss percentage & $-2.1798$ & 1.47 & \texttt{llm} \\
Options Delta & --- & 1.76 & \texttt{llm} \\
Collateral ratio & --- & 2.30 & \texttt{llm} \\
Funding rate cost (extended) & 1.0000 & 1.29 & \texttt{llm} \\
Staking reward for fixed lock-up & $-0.9246$ & 1.39 & \texttt{llm} \\
Expected Shortfall at 95\% & 1.0000 & 1.15 & \texttt{llm} \\
Expected Shortfall at 99\% & 1.0000 & 1.38 & \texttt{llm} \\
Optimal LP Position (Kelly) & 0.0000 & 2.59 & \texttt{nn\_fallback} \\
Concentrated liquidity position width & $-2.59\times10^{8}$ & 1.39 & \texttt{llm} \\
Portfolio VaR for two correlated & 1.0000 & 1.66 & \texttt{llm} \\
Position margin ratio & 1.0000 & 1.38 & \texttt{llm} \\
Concentrated liquidity position width (v2) & 1.0000 & 1.20 & \texttt{llm} \\
Spot price from AMM reserve & 0.9810 & 1.33 & \texttt{llm} \\
Black-Scholes Call Price & $-1.8315$ & 4.21 & \texttt{nn} \\
Black-Scholes Put Price & $-0.6802$ & 3.36 & \texttt{nn} \\
Component ES & $-0.2749$ & 1.16 & \texttt{llm} \\
Gamma of option & $-0.1060$ & 4.37 & \texttt{nn} \\
Vega of option & 0.6285 & 4.61 & \texttt{nn} \\
Multi-Collateral LTV & 1.0000 & 1.39 & \texttt{llm} \\
Correlated Portfolio VaR & 1.0000 & 1.44 & \texttt{llm} \\
Impermanent loss in constant product & 1.0000 & 1.26 & \texttt{llm} \\
Constant product formula (multivariate) & 1.0000 & 1.29 & \texttt{llm} \\
Convexity Adjustment & 0.9998 & 1.39 & \texttt{llm} \\
Liquidation price for leveraged long & --- & 1.32 & \texttt{llm} \\
Liquidation price for leveraged short & --- & 1.64 & \texttt{llm} \\
Maximum safe leverage & --- & 1.34 & \texttt{llm} \\
Required collateral & --- & 1.08 & \texttt{llm} \\
Forward price for derivative & 1.0000 & 1.92 & \texttt{llm} \\
Portfolio Expected Shortfall for correlated & --- & 1.34 & \texttt{llm} \\
Call option intrinsic & $-3.5798$ & 1.29 & \texttt{llm} \\
Put-call parity & 1.0000 & 1.36 & \texttt{llm} \\
Simple options moneyness & $-50.5674$ & 1.38 & \texttt{llm} \\
Uniswap V3 virtual & $-9.41\times10^{6}$ & 1.24 & \texttt{llm} \\
Theta of option & $-0.8034$ & 13.51 & \texttt{nn} \\
\end{longtable}
\endgroup


\subsection{Complete Test Suite Details}

\begin{table}[H]
\centering
\caption{Ground Truth Formulas with Exact Constants}
\label{tab:formulas_detailed}
\scriptsize
\begin{tabular}{lll}
\toprule
\textbf{Equation} & \textbf{Formula} & \textbf{Constants} \\
\midrule
Arrhenius & $k = A\exp(-E_a/(RT))$ & $A=10^{11}$, $E_a=80000$ J/mol, $R=8.314$ J/(mol$\cdot$K) \\
Henderson-Hasselbalch & $\text{pH} = \text{p}K_a + \log_{10}([A^-]/[\text{HA}])$ & $\text{p}K_a = 6.5$ \\
Rate Law & $r = k[A]^2[B]$ & $k = 0.5$ M$^{-2}$s$^{-1}$ \\
Allometric Scaling & $Y = aM^b$ & $a = 3.5$, $b = 0.75$ \\
Michaelis-Menten & $v = V_{\max}[S]/(K_m + [S])$ & $V_{\max} = 50$ $\mu$M/s, $K_m = 10$ $\mu$M \\
Logistic Growth & $dN/dt = rN(1-N/K)$ & $r = 0.3$ day$^{-1}$, $K = 1000$ \\
Kinetic Energy & $E = \frac{1}{2}mv^2$ & Dimensionless \\
Gravitational Force & $F = Gm_1m_2/r^2$ & $G = 6.674 \times 10^{-11}$ N$\cdot$m$^2$/kg$^2$ \\
Ideal Gas Law & $P = nRT/V$ & $R = 8.314$ J/(mol$\cdot$K) \\
Impermanent Loss & $\text{IL} = 2\sqrt{r}/(1+r) - 1$ & Dimensionless \\
Price Impact & $\Delta p = \Delta x/(x + \Delta x)$ & Dimensionless \\
Constant Product & $y = k/x$ & $k = 10^6$ \\
VaR 95\% & $\text{VaR} = P\sigma \cdot z_{0.95}$ & $z_{0.95} = 1.645$ \\
Liquidation Price & $p_{\text{liq}} = p_0(1 - 1/(L \cdot m))$ & $m = 0.8$ (maintenance margin) \\
Portfolio Variance & $\sigma_p = \sqrt{\sigma_1^2 + \sigma_2^2 + 2\rho\sigma_1\sigma_2}$ & Dimensionless \\
\bottomrule
\end{tabular}
\end{table}

% ============================================================================
% APPENDIX: NOISE SENSITIVITY ANALYSIS
% ============================================================================


\section{Noise Sensitivity Analysis}
\label{app:noise_sensitivity}

This section reports solve counts for the five noise levels present in the
repository's \texttt{protocol\_core} result files ($\sigma\leq1\%$). No results are available for higher noise levels, so no claim is made about a noise threshold.

\subsection{Protocol}

Five \texttt{protocol\_core} runs (one noiseless, four noisy; $200$ samples each)
cover $\sigma\in\{0,\,0.05,\,0.1,\,0.5,\,1\}\%$, read from each file's
\texttt{protocol.noise\_level}, with 30 equations per level (150 records per
system). A record counts as solved if $\Rsq\geq$ the threshold recorded in its
own file's protocol block (\texttt{protocol.threshold}); the
\texttt{success} flag is not used (M3 sets it to true whenever no exception is
raised).

\subsection{Results}

\begin{table}[H]
\centering
\caption{Solved counts per noise level ($\Rsq\geq$ the run's own recorded
  threshold; $n=30$ per cell). From the five \texttt{protocol\_core}
  files; columns give $\sigma$ in percent. \textbf{Thresholds differ per run}
  ($0.999999$, $0.95$, $0.95$, $0.995$, $0.99$ for $\sigma=0$, $0.05$,
  $0.1$, $0.5$, $1\%$; these are the thresholds written into each
  \texttt{protocol\_core} file, and differ from the $0.995$ that
  \Cref{tab:rr_noise} applies at $\sigma=0.05\%$ and $0.1\%$), so counts are \emph{not} comparable across noise
  levels, and the pooled column adds counts taken under different thresholds.
  \textbf{Source of the $\sigma=0$ column:} all six cells come from the noiseless run of the noise sweep, \texttt{protocol\_core\_noiseless\_20260812\_170344.json}; no cell is taken from \texttt{benchmark\_results.json} (see the tablenote on \Cref{tab:overall}).}
\label{tab:noise_sensitivity}
\small
\setlength{\tabcolsep}{4pt}
\begin{tabular}{@{}l c c c c c r@{}}
\toprule
\textbf{Method, $\sigma$ (\%)} & \textbf{0} & \textbf{0.05} & \textbf{0.1} & \textbf{0.5} & \textbf{1} & \textbf{Pooled}\\
\midrule
\EHD{} (M3)            & 23/30 & 30/30 & 30/30 & 30/30 & 30/30 & 143/150 (95.3\%)\\
HybridDiscovery v50\_2 & 21/30 & 25/30 & 26/30 & 23/30 & 26/30 & 121/150 (80.7\%)\\
\HSL{} (M4)            & 30/30 & 30/30 & 30/30 & 30/30 & 30/30 & 150/150 (100.0\%)\\
ImprovedNN             &  0/30 & 29/30 & 29/30 & 27/30 & 29/30 & 114/150 (76.0\%)\\
PureLLM                & 30/30 & 29/30 & 29/30 & 30/30 & 30/30 & 148/150 (98.7\%)\\
SymbolicEngineWithLLM  & 28/30 & 30/30 & 30/30 & 29/30 & 29/30 & 146/150 (97.3\%)\\
\bottomrule
\end{tabular}
\end{table}

\begin{table}[H]
\centering
\caption{Solved counts at a single fixed threshold, $\Rsq\geq0.99$, for every
  run ($n=30$ per cell). Because the threshold is the same at all noise levels,
  counts are comparable across levels and the pooled column is valid.
  Columns give $\sigma$ in percent. The $\sigma=0$ column is read from \texttt{protocol\_core\_noiseless\_20260812\_170344.json}, as in \Cref{tab:noise_sensitivity}.}
\label{tab:noise_sensitivity_fixed}
\small
\setlength{\tabcolsep}{4pt}
\begin{tabular}{@{}l c c c c c r@{}}
\toprule
\textbf{Method, $\sigma$ (\%)} & \textbf{0} & \textbf{0.05} & \textbf{0.1} & \textbf{0.5} & \textbf{1} & \textbf{Pooled}\\
\midrule
\EHD{} (M3)            & 30/30 & 30/30 & 30/30 & 30/30 & 30/30 & 150/150 (100.0\%)\\
HybridDiscovery v50\_2 & 30/30 & 24/30 & 26/30 & 24/30 & 26/30 & 130/150 (86.7\%)\\
\HSL{} (M4)            & 30/30 & 30/30 & 30/30 & 30/30 & 30/30 & 150/150 (100.0\%)\\
ImprovedNN             & 29/30 & 27/30 & 28/30 & 29/30 & 29/30 & 142/150 (94.7\%)\\
PureLLM                & 30/30 & 29/30 & 29/30 & 30/30 & 30/30 & 148/150 (98.7\%)\\
SymbolicEngineWithLLM  & 30/30 & 30/30 & 30/30 & 30/30 & 29/30 & 149/150 (99.3\%)\\
\bottomrule
\end{tabular}
\end{table}

\subsection{Key Findings}

\paragraph{Ranking.} At the fixed $\Rsq\geq0.99$ threshold, M3 and M4 both solve
150/150, ahead of SymbolicEngineWithLLM (149), PureLLM (148), ImprovedNN (142)
and HybridDiscovery v50\_2 (130). M3's apparent shortfall in the per-run table
(23/30 at $\sigma=0$) comes from the strict $0.999999$ threshold used for that
run, not from a loss of accuracy. At that strict threshold M3 solves 23/30
noiseless equations while M4 solves 30/30, so M4 is the only system that is exact
under the strict criterion. At $0.99$ the two systems cannot be separated.

\paragraph{Threshold confound.} The per-run table mixes thresholds from $0.95$ to
$0.999999$. For example ImprovedNN scores 0/30 at $\sigma=0$ under $0.999999$ but
29/30 under $0.99$; the jump in the per-run table is the threshold, not the
data. Cross-level comparisons should use the fixed-threshold table. The value
$0.99$ is permissive, so it does not measure exact recovery.

\paragraph{Scope.} Only $\sigma\leq1\%$ is covered, and at
$\Rsq\geq0.99$ essentially all systems saturate, so these data do not show where
noise starts to hurt. No failure-mode, domain-variability, or noise-threshold claim is supported by
these data.

%% ============================================================
%% Material relocated from the main paper (50-page reduction)
%% ============================================================
\section{Material Relocated from the Main Paper (Benchmarks)}\label{sec:relocated_bench}
This section holds tables, a figure and case definitions removed from the main paper to keep it within 50 pages. Text is reproduced verbatim; references to sections and tables of the main paper are resolved through the main paper's auxiliary file.

\subsection{Protocol-30: per-equation results (random 80/20 and PCA 40/60 splits)}
{\scriptsize\setlength{\tabcolsep}{2.5pt}
\begin{longtable}{p{1.9cm}p{4.3cm}cccccc}
\caption{Random 80/20 split, run dated 2026-08-12, all 30 equations $\times$ six methods (method-reported $\Rsq$, \texttt{benchmark\_results.json}). \textbf{Bold}: pass ($\Rsq \geq 0.999999$). PL = PureLLM Baseline, NN = ImprovedNN, EH = EnhancedHybridSystemDeFi, LN = HybridSystemLLMNN all-domains, SE = SymbolicEngineWithLLM, HD = HybridDiscoverySystem v50\_2 (\textsc{HypatiaX}). The PureLLM entry for Curie's law is a recorded evaluation failure (no valid formula), shown as $0.000000$.}
\label{tab:randomsplit}\\
\toprule
Domain & Equation & PL & NN & EH & LN & SE & HD \\
\midrule
\endfirsthead
\toprule
Domain & Equation & PL & NN & EH & LN & SE & HD \\
\midrule
\endhead
Biology & Michaelis--Menten kinetics & \textbf{1.000000} & 0.999997 & 0.999988 & \textbf{1.000000} & \textbf{1.000000} & \textbf{1.000000} \\
Biology & Logistic growth rate & \textbf{1.000000} & 0.999859 & 0.999997 & \textbf{1.000000} & \textbf{1.000000} & \textbf{1.000000} \\
Biology & Allometric scaling law & \textbf{1.000000} & 0.999995 & 0.999993 & \textbf{1.000000} & \textbf{1.000000} & \textbf{1.000000} \\
Chemistry & Arrhenius rate constant & \textbf{1.000000} & 0.999930 & 0.999994 & \textbf{1.000000} & 0.999921 & 0.999594 \\
Chemistry & Henderson--Hasselbalch buffer pH & \textbf{1.000000} & 0.999985 & \textbf{1.000000} & \textbf{1.000000} & \textbf{1.000000} & \textbf{1.000000} \\
Electrochem. & Nernst equation & \textbf{1.000000} & 0.999979 & \textbf{1.000000} & \textbf{1.000000} & 0.997725 & 0.999205 \\
Electromagn. & Clausius--Mossotti & \textbf{1.000000} & 0.999973 & \textbf{1.000000} & \textbf{1.000000} & \textbf{1.000000} & \textbf{1.000000} \\
Electromagn. & Dielectric polarisation & \textbf{1.000000} & 0.999242 & \textbf{1.000000} & \textbf{1.000000} & \textbf{1.000000} & 0.995528 \\
Electromagn. & Lorentz force & \textbf{1.000000} & 0.999975 & \textbf{1.000000} & \textbf{1.000000} & \textbf{1.000000} & \textbf{1.000000} \\
Electromagn. & Ohm's law & \textbf{1.000000} & 0.999587 & \textbf{1.000000} & \textbf{1.000000} & \textbf{1.000000} & \textbf{1.000000} \\
Electromagn. & Energy in a capacitor & \textbf{1.000000} & 0.999806 & \textbf{1.000000} & \textbf{1.000000} & \textbf{1.000000} & 0.999964 \\
Electrostat. & Coulomb force (1D, simplified) & \textbf{1.000000} & 0.996355 & \textbf{1.000000} & \textbf{1.000000} & \textbf{1.000000} & \textbf{1.000000} \\
Electrostat. & Coulomb's law & \textbf{1.000000} & 0.999939 & \textbf{1.000000} & \textbf{1.000000} & \textbf{1.000000} & 0.996597 \\
Magnetism & Curie's law & 0.000000 & 0.999953 & \textbf{1.000000} & \textbf{1.000000} & \textbf{1.000000} & \textbf{1.000000} \\
Mechanics & Newton's gravitational force & \textbf{1.000000} & 0.993960 & \textbf{1.000000} & \textbf{1.000000} & \textbf{1.000000} & \textbf{1.000000} \\
Mechanics & Kinetic energy & \textbf{1.000000} & 0.999855 & \textbf{1.000000} & \textbf{1.000000} & \textbf{1.000000} & 0.999878 \\
Mechanics & Reduced mass & \textbf{1.000000} & 0.999995 & \textbf{1.000000} & \textbf{1.000000} & \textbf{1.000000} & \textbf{1.000000} \\
Mechanics & Total mechanical energy & \textbf{1.000000} & 0.999937 & \textbf{1.000000} & \textbf{1.000000} & \textbf{1.000000} & \textbf{1.000000} \\
Optics & Snell's law & \textbf{1.000000} & 0.999987 & \textbf{1.000000} & \textbf{1.000000} & \textbf{1.000000} & \textbf{1.000000} \\
Optics & Double-slit interference intensity & \textbf{1.000000} & 0.999902 & \textbf{1.000000} & \textbf{1.000000} & \textbf{1.000000} & \textbf{1.000000} \\
Probability & Gaussian probability density & \textbf{1.000000} & 0.999770 & \textbf{1.000000} & \textbf{1.000000} & \textbf{1.000000} & 0.999765 \\
Quantum & Photon energy & \textbf{1.000000} & 0.999998 & 0.999990 & \textbf{1.000000} & \textbf{1.000000} & \textbf{1.000000} \\
Quantum & Zeeman energy & \textbf{1.000000} & 0.999910 & \textbf{1.000000} & \textbf{1.000000} & \textbf{1.000000} & \textbf{1.000000} \\
Quantum & Bose--Einstein occupation number & \textbf{1.000000} & 0.999960 & 0.999979 & \textbf{1.000000} & \textbf{1.000000} & 0.994634 \\
Quantum & Fermi--Dirac occupation number & \textbf{1.000000} & 0.998318 & \textbf{1.000000} & \textbf{1.000000} & 0.998390 & 0.999037 \\
Quantum & Rabi frequency & \textbf{1.000000} & 0.999946 & \textbf{1.000000} & \textbf{1.000000} & \textbf{1.000000} & \textbf{1.000000} \\
Thermodyn. & Planck blackbody radiance & \textbf{1.000000} & 0.999932 & 0.999996 & \textbf{1.000000} & \textbf{1.000000} & \textbf{1.000000} \\
Thermodyn. & Fourier's law & \textbf{1.000000} & 0.993764 & \textbf{1.000000} & \textbf{1.000000} & \textbf{1.000000} & \textbf{1.000000} \\
Thermodyn. & Stefan--Boltzmann law & \textbf{1.000000} & 0.999975 & \textbf{1.000000} & \textbf{1.000000} & \textbf{1.000000} & \textbf{1.000000} \\
Thermodyn. & Ideal gas law & \textbf{1.000000} & 0.999911 & \textbf{1.000000} & \textbf{1.000000} & \textbf{1.000000} & \textbf{1.000000} \\
\bottomrule
\end{longtable}}


{\scriptsize\setlength{\tabcolsep}{2.5pt}
\begin{longtable}{p{1.9cm}p{4.3cm}cccccc}
\caption{PCA-directed 40/60 split, run dated 2026-08-12, all 30 equations $\times$ six methods (method-reported $\Rsq$, \texttt{benchmark\_results\_pca\_4060.json}). \textbf{Bold}: pass ($\Rsq \geq 0.999999$). Column abbreviations as in Table~\ref{tab:randomsplit}. The PureLLM entry for the capacitor-energy equation is a recorded evaluation failure, shown as $0.000000$.}
\label{tab:pcasplit}\\
\toprule
Domain & Equation & PL & NN & EH & LN & SE & HD \\
\midrule
\endfirsthead
\toprule
Domain & Equation & PL & NN & EH & LN & SE & HD \\
\midrule
\endhead
Biology & Michaelis--Menten kinetics & \textbf{1.000000} & 0.993994 & 0.999997 & \textbf{1.000000} & \textbf{1.000000} & \textbf{1.000000} \\
Biology & Logistic growth rate & \textbf{1.000000} & 0.994709 & 0.999998 & \textbf{1.000000} & \textbf{1.000000} & \textbf{1.000000} \\
Biology & Allometric scaling law & \textbf{1.000000} & 0.936331 & 0.999998 & \textbf{1.000000} & \textbf{1.000000} & \textbf{1.000000} \\
Chemistry & Arrhenius rate constant & \textbf{1.000000} & -0.044405 & 0.999996 & \textbf{1.000000} & 0.994936 & 0.997920 \\
Chemistry & Henderson--Hasselbalch buffer pH & \textbf{1.000000} & 0.972446 & \textbf{1.000000} & \textbf{1.000000} & 0.999283 & 0.999951 \\
Electrochem. & Nernst equation & \textbf{1.000000} & 0.990615 & \textbf{1.000000} & \textbf{1.000000} & 0.998771 & 0.999983 \\
Electromagn. & Clausius--Mossotti & \textbf{1.000000} & 0.909288 & \textbf{1.000000} & \textbf{1.000000} & \textbf{1.000000} & \textbf{1.000000} \\
Electromagn. & Dielectric polarisation & \textbf{1.000000} & 0.710588 & \textbf{1.000000} & \textbf{1.000000} & \textbf{1.000000} & 0.995048 \\
Electromagn. & Lorentz force & \textbf{1.000000} & 0.800979 & \textbf{1.000000} & \textbf{1.000000} & \textbf{1.000000} & \textbf{1.000000} \\
Electromagn. & Ohm's law & \textbf{1.000000} & 0.960263 & \textbf{1.000000} & \textbf{1.000000} & \textbf{1.000000} & \textbf{1.000000} \\
Electromagn. & Energy in a capacitor & 0.000000 & 0.855165 & \textbf{1.000000} & \textbf{1.000000} & \textbf{1.000000} & \textbf{1.000000} \\
Electrostat. & Coulomb force (1D, simplified) & \textbf{1.000000} & 0.999966 & \textbf{1.000000} & \textbf{1.000000} & \textbf{1.000000} & \textbf{1.000000} \\
Electrostat. & Coulomb's law & \textbf{1.000000} & 0.999862 & \textbf{1.000000} & \textbf{1.000000} & \textbf{1.000000} & 0.999615 \\
Magnetism & Curie's law & \textbf{1.000000} & 0.998748 & \textbf{1.000000} & \textbf{1.000000} & \textbf{1.000000} & \textbf{1.000000} \\
Mechanics & Newton's gravitational force & \textbf{1.000000} & 0.999990 & \textbf{1.000000} & \textbf{1.000000} & \textbf{1.000000} & \textbf{1.000000} \\
Mechanics & Kinetic energy & \textbf{1.000000} & 0.112559 & \textbf{1.000000} & \textbf{1.000000} & \textbf{1.000000} & 0.995804 \\
Mechanics & Reduced mass & \textbf{1.000000} & 0.925415 & \textbf{1.000000} & \textbf{1.000000} & \textbf{1.000000} & \textbf{1.000000} \\
Mechanics & Total mechanical energy & \textbf{1.000000} & 0.939856 & \textbf{1.000000} & \textbf{1.000000} & \textbf{1.000000} & \textbf{1.000000} \\
Optics & Snell's law & \textbf{1.000000} & 0.902944 & \textbf{1.000000} & \textbf{1.000000} & \textbf{1.000000} & \textbf{1.000000} \\
Optics & Double-slit interference intensity & \textbf{1.000000} & 0.733607 & \textbf{1.000000} & \textbf{1.000000} & \textbf{1.000000} & \textbf{1.000000} \\
Probability & Gaussian probability density & \textbf{1.000000} & 0.489964 & \textbf{1.000000} & \textbf{1.000000} & 0.997249 & 0.997945 \\
Quantum & Photon energy & \textbf{1.000000} & \textbf{1.000000} & 0.999991 & \textbf{1.000000} & \textbf{1.000000} & \textbf{1.000000} \\
Quantum & Zeeman energy & \textbf{1.000000} & 0.786961 & \textbf{1.000000} & \textbf{1.000000} & \textbf{1.000000} & \textbf{1.000000} \\
Quantum & Bose--Einstein occupation number & \textbf{1.000000} & 0.999995 & 0.999890 & \textbf{1.000000} & \textbf{1.000000} & 0.998929 \\
Quantum & Fermi--Dirac occupation number & \textbf{1.000000} & 0.927491 & \textbf{1.000000} & \textbf{1.000000} & 0.511948 & 0.647311 \\
Quantum & Rabi frequency & \textbf{1.000000} & 0.945751 & \textbf{1.000000} & \textbf{1.000000} & \textbf{1.000000} & \textbf{1.000000} \\
Thermodyn. & Planck blackbody radiance & \textbf{1.000000} & -0.104423 & 0.999997 & \textbf{1.000000} & 0.999946 & 0.999969 \\
Thermodyn. & Fourier's law & \textbf{1.000000} & 0.835533 & \textbf{1.000000} & \textbf{1.000000} & 0.999853 & 0.948649 \\
Thermodyn. & Stefan--Boltzmann law & \textbf{1.000000} & \textbf{1.000000} & \textbf{1.000000} & \textbf{1.000000} & \textbf{1.000000} & 0.999804 \\
Thermodyn. & Ideal gas law & \textbf{1.000000} & \textbf{1.000000} & \textbf{1.000000} & \textbf{1.000000} & \textbf{1.000000} & \textbf{1.000000} \\
\bottomrule
\end{longtable}}

\subsection{Five-system comparison, Experiment 1: interpolation performance}
\begin{table}[t]
\centering
\caption{Five-System Comparison -- Performance (Core-15, interpolation).
  95\% CI computed from raw per-equation train $R^2$/RMSE at generation time.}
\label{tab:five_systems_performance}
\begin{tabular}{lrrrrr}
\toprule
\textbf{System} & \textbf{n} & \textbf{Train $R^2$ Mean}
  & \textbf{Train $R^2$ 95\% CI} & \textbf{Train RMSE Mean}
  & \textbf{Design Focus} \\
\midrule
Pure LLM & 15 & -3.077 & [-10.705, 4.550] & 342.8347 & Recognition \\
Neural Network & 15 & 0.982 & [0.957, 1.007] & 1999.5749 & Baseline \\
System 3 LLM+Fallback & 15 & \textbf{1.000} & [0.999, 1.000] & \textbf{2.9626} & Robustness \\
System 2 Symbolic & 15 & 0.998 & [0.998, 0.998] & 236.1183 & Validation \\
Hybrid v50\_2 & 15 & 0.997 & [0.995, 0.998] & 236.1300 & Extrapolation \\
\bottomrule
\end{tabular}
\end{table}

\subsection{Nguyen-12: recoveries per equation across five seeds}
\begin{table}[h]
\centering
\caption{Nguyen-12, recoveries per equation across the five seeds (42, 99, 123, 777, 2024) at $\Rsq\ge0.9999$. $H$ = HypatiaX, $P$ = PySR-only.}
\label{tab:nguyen12_byeq}
\begin{tabular}{lrr}
\toprule
Eq. & HypatiaX ($H$) & PySR-only ($P$) \\
\midrule
N-1 & 5/5 & 5/5 \\
N-2 & 5/5 & 5/5 \\
N-3 & 5/5 & 4/5 \\
N-4 & 5/5 & 2/5 \\
N-5 & 2/5 & 1/5 \\
N-6 & 5/5 & 5/5 \\
N-7 & 5/5 & 3/5 \\
N-8 & 5/5 & 5/5 \\
N-9 & 4/5 & 4/5 \\
N-10 & 5/5 & 5/5 \\
N-11 & 5/5 & 5/5 \\
N-12 & 0/5 & 5/5 \\
\bottomrule
\end{tabular}
\end{table}

\subsection{Core-15 per-equation $R^2$ heatmap by extrapolation regime}
\begin{figure}[htbp]
\centering
\includegraphics[width=\linewidth]{fig09_r2_heatmap_regimes}
\caption{$\Rsq$ heatmap by equation (rows) and extrapolation regime
(columns: Train, Near $1.2\times$, Medium, Far $5\times$), Core-15
benchmark, on a shared colour scale saturating outside $[-1.5, 1.0]$.
\textbf{This figure is generated from the values in the Core-15 ablation table in the main paper} and the two
agree. For five of fifteen equations HypatiaX Hybrid returns \texttt{nan} (Henderson-Hasselbalch,
Logistic Growth, Michaelis-Menten, Rate Law; grey hatched) or $-\infty$ (Gravitational Force)
beyond the Train column --- missing or crashed evaluations, not $\Rsq$ scores of zero, and not
to be confused with a poor fit. Constant Product completes with finite but very large
negative values ($-173$, $-246$, $-4.5\times10^{4}$ at Near, Medium, Far). PySR-only
reports a value in every cell except Henderson-Hasselbalch Medium ($-\infty$); its most severe
finite far values are Michaelis-Menten ($-808$), Impermanent Loss ($-63$), Henderson-Hasselbalch
($-12$) and Arrhenius ($-12.55$). On Arrhenius, HypatiaX Hybrid's completed Far value ($+0.90$) is a \emph{success}
(\S\ref{sec:arrhenius_failure}). Apart from Constant Product (Far $=-4.5\times10^{4}$), HypatiaX Hybrid's most severe
completed value is Impermanent Loss (Far $=-8259$).}
\label{fig:r2_heatmap_raw_s}
\end{figure}

\subsection{Full Benchmark Case Definitions}
\subsubsection{Automated Market Makers}\label{subsec:automated-market-makers}
\label{sec:bench_amm}

\begin{itemize}
  \item \textbf{Constant product formula}: $x \cdot y = k$.
    Tests basic AMM invariant preservation.
  \item \textbf{Spot price from AMM}: $P = y/x$.
    Algebraic; should be trivially recovered.
  \item \textbf{Spot price from AMM reserve}: $P = \Delta y / \Delta x$.
    Marginal price from reserve changes.
  \item \textbf{AMM output amount}: $\Delta y = y - k/(x + \Delta x)$.
    Tests multi-step algebraic composition.
  \item \textbf{Price slippage percentage}: $\mathrm{slip} = (P_{\mathrm{exec}} - P_{\mathrm{spot}})/P_{\mathrm{spot}}$.
  \item \textbf{Constant Product Price Impact}: quadratic impact formula.
  \item \textbf{LP share percentage}: $s = \Delta L / L_{\mathrm{total}}$.
  \item \textbf{LP fee earnings}: $e = s \cdot F \cdot V$, where $F$ is fee tier and $V$ volume.
  \item \textbf{Impermanent loss percentage}: $\mathrm{IL} = 2\sqrt{r}/(1+r) - 1$,
    where $r$ is the price ratio at withdrawal vs.\ deposit.
  \item \textbf{Impermanent loss in constant product}: related quantity for the
    constant-product setting specifically.
  \item \textbf{Impermanent loss breakeven fee rate}: the fee rate $F$ at which
    LP fees exactly offset impermanent loss.
  \item \textbf{Uniswap V3 virtual}: virtual liquidity quantity for
    concentrated positions.
  \item \textbf{Concentrated liquidity position width}: tick range formula.
  \item \textbf{Capital efficiency}: ratio of effective liquidity to deposited capital.
\end{itemize}

\subsubsection{Lending Protocols}\label{subsec:lending-protocols}
\label{sec:bench_lending}

\begin{itemize}
  \item \textbf{Loan-to-Value}: $\mathrm{LTV} = D/C$.
  \item \textbf{Health factor}: $h = C_{\mathrm{adj}} / D$.
  \item \textbf{Collateral ratio}: $\mathrm{CR} = C/D$.
  \item \textbf{Reserve ratio}: $\rho = R / (R + L)$.
  \item \textbf{Utilization rate of DeFi}: $u = L / (L + R)$.
  \item \textbf{Borrow APY from utilization}: kinked interest rate model.
  \item \textbf{Supply APY from borrow}: $r_s = r_b \cdot u$.
  \item \textbf{Partial liquidation amount}: amount of collateral seized.
  \item \textbf{Liquidation Price Long} and \textbf{Liquidation Price Short}:
    prices at which a position is liquidated.
  \item \textbf{Leveraged position notional}: $N = C \cdot L$, where $L$ is leverage.
  \item \textbf{Effective Leverage}: $L_{\mathrm{eff}} = N/C$.
  \item \textbf{Maximum safe leverage}: function of maintenance margin rate.
  \item \textbf{Required collateral}: collateral needed for a given position.
  \item \textbf{Cross-margin available balance}: net balance accounting for all positions.
  \item \textbf{Multi-Collateral LTV}: weighted LTV across multiple collateral types.
  \item \textbf{Funding-Adjusted Liquidation Price (Long)} and
    \textbf{Funding-Adjusted Liquidation Price (Short)}: extended formulas accounting for funding rates.
\end{itemize}

\subsubsection{Derivatives Pricing}\label{subsec:derivatives-pricing}
\label{sec:bench_derivatives}

\begin{itemize}
  \item \textbf{Black--Scholes Call Price}:
    $C = S\Phi(d_1) - Ke^{-rT}\Phi(d_2)$,
    with $d_{1,2} = [\ln(S/K) + (r \pm \tfrac{1}{2}\sigma^2)T] / (\sigma\sqrt{T})$.
  \item \textbf{Black--Scholes Put Price}: $P = Ke^{-rT}\Phi(-d_2) - S\Phi(-d_1)$.
  \item \textbf{Delta}: $\Delta = \Phi(d_1)$ (call).
  \item \textbf{Gamma}: $\Gamma = \phi(d_1)/(S\sigma\sqrt{T})$.
  \item \textbf{Theta}: $\Theta = -S\phi(d_1)\sigma/(2\sqrt{T}) - rKe^{-rT}\Phi(d_2)$.
  \item \textbf{Vega}: $\mathcal{V} = S\phi(d_1)\sqrt{T}$.
  \item \textbf{Put-call parity}: $C - P = S - Ke^{-rT}$.
  \item \textbf{Forward price for derivative}: $F = Se^{rT}$.
  \item \textbf{Call option intrinsic}: $\max(S - K, 0)$.
  \item \textbf{Simple options moneyness}: $(S - K)/K$.
  \item \textbf{Convexity Adjustment}: bond convexity formula.
\end{itemize}

\subsubsection{Risk and Portfolio Metrics}\label{subsec:risk-and-portfolio-metrics}
\label{sec:bench_risk}

\begin{itemize}
  \item \textbf{Value at Risk at 95\,\% and 99\,\%}: parametric VaR formulas.
  \item \textbf{Expected Shortfall at 95\,\% and 99\,\%}: conditional VaR formulas.
  \item \textbf{Multi-day Value at Risk}: $\mathrm{VaR}_{T} = \mathrm{VaR}_1 \cdot \sqrt{T}$.
  \item \textbf{ES scaling for multi-day}: analogous scaling for ES.
  \item \textbf{Incremental VaR}: marginal contribution of one asset to portfolio VaR.
  \item \textbf{Portfolio VaR for two correlated assets}:
    $\mathrm{VaR}_P = \sqrt{w_1^2 \sigma_1^2 + w_2^2 \sigma_2^2 + 2w_1 w_2 \rho\sigma_1\sigma_2}$.
  \item \textbf{Correlated Portfolio VaR}: multivariate extension.
  \item \textbf{Component ES}: ES contribution decomposed by asset.
  \item \textbf{Portfolio Expected Shortfall for correlated assets}: the
    most complex risk case; requires integration over a multivariate Gaussian tail.
  \item \textbf{Annualised Portfolio tracking error}: $\mathrm{TE} = \sigma_{\varepsilon}\sqrt{252}$.
  \item \textbf{Portfolio Sharpe Ratio}: $(r_p - r_f)/\sigma_p$.
  \item \textbf{Information ratio}: $\mathrm{IR} = \alpha / \mathrm{TE}$.
  \item \textbf{Position margin ratio}: margin / position notional.
\end{itemize}

\subsubsection{Staking and Protocol Mechanisms}\label{subsec:staking-and-protocol-mechanisms}
\label{sec:bench_staking}

\begin{itemize}
  \item \textbf{Simple Staking APY}: $\mathrm{APY} = r$.
  \item \textbf{Compounding Staking Returns}: $(1 + r/n)^{nT} - 1$.
  \item \textbf{Staking reward for fixed lock-up}: $R = P \cdot r \cdot T$.
  \item \textbf{Validator commission adjusted}: post-commission yield.
  \item \textbf{Slashing penalty}: fractional penalty applied to staked balance.
  \item \textbf{Protocol reserve accumulation}: cumulative reserve growth formula.
  \item \textbf{Funding rate cost}: $C = N \cdot f \cdot h$, where $f$ is the funding
    rate and $h$ is the holding period in funding intervals.
  \item \textbf{APY calculation}: annualised yield from periodic rates.
\end{itemize}


\end{document}
